\documentclass[11pt]{article}
\usepackage[margin=1in]{geometry}
\usepackage[T1]{fontenc}
\usepackage[utf8]{inputenc}
\usepackage[french]{babel}
\usepackage{amsmath,amssymb,amsthm}
\usepackage[colorlinks=true,linkcolor=blue,urlcolor=blue]{hyperref}
\newtheorem{problem}{Problème}
\newenvironment{refs}{\par\small\noindent\emph{Références :}\begin{itemize}\setlength\itemsep{0pt}}{\end{itemize}\normalsize}
\title{Hors de la Caverne \\ \Large Algèbre linéaire}
\author{}
\date{}
\begin{document}
\maketitle
\noindent\emph{Des problèmes, pas de réponses.} Chaque problème ci-dessous est posé
sans solution. Les références indiquent où trouver les connaissances nécessaires.
\medskip


\section*{Lycée}

\begin{problem}[{Deux droites, une réponse --- le plus souvent --- Lycée}]
\textbf{LIN-01.} \textbf{Problème.} Considérons le système de deux équations
$ax + by = e$, $cx + dy = f$, où $a, b, c, d, e, f$ sont des nombres réels donnés.
\begin{enumerate}
\item[(a)] Démontrer que le système possède exactement une solution $(x, y)$ si et seulement si
$ad - bc \ne 0$.
\item[(b)] Lorsque $ad - bc = 0$, déterminer --- avec démonstration --- précisément quand le système
n'a aucune solution et quand il en a une infinité.
\item[(c)] Interpréter géométriquement chaque cas comme un
énoncé portant sur deux droites du plan.
\end{enumerate}

\end{problem}

\noindent Résoudre des systèmes d'équations est plus ancien que l'algèbre elle-même : le
chapitre 8 du classique chinois \textit{Les Neuf Chapitres sur l'art mathématique} (compilé au
Ier siècle de notre ère) résout de tels systèmes en manipulant le tableau des
coefficients --- le pivot de Gauss, dix-huit siècles avant Gauss, qui le réinventa en 1809 pour
suivre l'astéroïde Cérès. La quantité $ad - bc$ qui fait sa première apparition en (i) est le
déterminant, dont cette page suit la longue carrière.

\begin{refs}
\item Contexte : Système d'équations linéaires --- Wikipédia (\url{https://fr.wikipedia.org/wiki/Syst%C3%A8me_d%27%C3%A9quations_lin%C3%A9aires})
\item Contexte : Les Neuf Chapitres sur l'art mathématique --- Wikipédia (\url{https://fr.wikipedia.org/wiki/Les_Neuf_Chapitres_sur_l%27art_math%C3%A9matique})
\item Lecture : G. Strang, \textit{Introduction to Linear Algebra}, ch. 1--2
\end{refs}

\bigskip

\begin{problem}[{Les rotations vues comme des matrices --- Lycée}]
\textbf{LIN-02.} \textbf{Problème.}
\begin{enumerate}
\item[(a)] Montrer que la rotation du plan autour de l'origine d'un angle
$\theta$ envoie le point $(x, y)$ sur
$(x \cos\theta - y \sin\theta, \; x \sin\theta + y \cos\theta)$.
\item[(b)] Démontrer qu'effectuer
la rotation d'angle $\alpha$ puis la rotation d'angle $\beta$ donne la rotation d'angle
$\alpha + \beta$, et en déduire les formules d'addition du sinus et du cosinus.
\item[(c)] Une réflexion
du plan par rapport à une droite passant par l'origine déplace elle aussi les points linéairement.
Trouver sa formule en fonction de l'angle de la droite, et démontrer que la composée de deux
réflexions est une rotation.
\end{enumerate}

\end{problem}

\noindent Voici la porte par laquelle la géométrie entre dans l'algèbre : les transformations
deviennent des tableaux de nombres, et la composition devient une multiplication. La notation
matricielle et la règle de multiplication ont été fixées par Cayley dans son \textit{Memoir on the Theory
of Matrices} de 1858, précisément pour rendre automatique une comptabilité comme celle de (ii).
La surprise agréable que les identités trigonométriques soient des théorèmes \textit{sur les matrices}
est, pour bien des étudiants, le premier indice qu'une notation peut penser à votre place.

\begin{refs}
\item Contexte : Matrice de rotation --- Wikipédia (\url{https://fr.wikipedia.org/wiki/Matrice_de_rotation})
\item Contexte : Matrice d'une application linéaire --- Wikipédia (\url{https://fr.wikipedia.org/wiki/Matrice_d%27une_application_lin%C3%A9aire})
\item Lecture : G. Strang, \textit{Introduction to Linear Algebra}, ch. 2
\end{refs}

\bigskip

\begin{problem}[{Le déterminant est une aire --- Lycée}]
\textbf{LIN-03.} \textbf{Problème.} Soient $u = (a, b)$ et $v = (c, d)$ deux vecteurs du
plan.
\begin{enumerate}
\item[(a)] Démontrer que l'aire du parallélogramme de côtés $u$ et $v$ vaut
$|ad - bc|$.
\item[(b)] Montrer que le signe de $ad - bc$ enregistre l'orientation : il est positif
exactement lorsque le plus court tour de $u$ vers $v$ se fait dans le sens direct.
\item[(c)] Utiliser (i) et (ii)
pour démontrer la \og{} formule du lacet \fg{} : l'aire d'un polygone de sommets
$(x_1, y_1), \dots, (x_n, y_n)$, énumérés dans le sens direct, vaut
$\frac{1}{2} \sum_{i=1}^{n} (x_i y_{i+1} - x_{i+1} y_i)$, les indices étant pris modulo $n$.
\end{enumerate}

\end{problem}

\noindent Les déterminants précèdent les matrices : Seki Takakazu au Japon et Leibniz en
Europe les rencontrèrent tous deux en 1683, comme conditions de compatibilité d'un système
d'équations. Que la même expression mesure une aire --- et, en dimension supérieure, un volume --- est
l'âme géométrique de la théorie, développée par Lagrange et Cauchy. La formule du lacet, outil de
base des arpenteurs, est habituellement attribuée à Meister (1769) et à Gauss ; c'est le déterminant
faisant honnêtement du travail de plein air.

\begin{refs}
\item Contexte : Déterminant (mathématiques) --- Wikipédia (\url{https://fr.wikipedia.org/wiki/D%C3%A9terminant_%28math%C3%A9matiques%29})
\item Contexte : Formule du lacet --- Wikipédia (en anglais) (\url{https://en.wikipedia.org/wiki/Shoelace_formula})
\item Lecture : G. Strang, \textit{Introduction to Linear Algebra}, ch. 5
\end{refs}

\bigskip

\begin{problem}[{Trace et déterminant cachés dans un polynôme --- Lycée, short}]
\textbf{LIN-22.} \textbf{Problème.} Pour la matrice
$A = \begin{pmatrix} a & b \\ c & d \end{pmatrix}$, démontrer que
\[ \det(A - \lambda I) = \lambda^2 - (a + d)\lambda + (ad - bc). \]
En déduire que si $A$ possède deux valeurs propres réelles, leur somme est la trace $a + d$ et
leur produit le déterminant $ad - bc$.
\end{problem}

\noindent Les deux nombres que tout le monde apprend à calculer pour une matrice
$2 \times 2$ se révèlent être les seuls coefficients de son polynôme caractéristique --- ce sont donc
exactement les données qu'une valeur propre peut voir. Cayley et Sylvester, qui baptisèrent la trace
et le déterminant dans les années 1850, cherchaient précisément de telles quantités : des expressions
des coefficients qui ne changent pas lorsque la matrice est réécrite dans une autre base. En
dimension $n$, le même polynôme a $n$ coefficients, et ce sont les fonctions symétriques
élémentaires des valeurs propres.

\begin{refs}
\item Contexte : Polynôme caractéristique --- Wikipédia (\url{https://fr.wikipedia.org/wiki/Polyn%C3%B4me_caract%C3%A9ristique})
\item Contexte : Trace (algèbre) --- Wikipédia (\url{https://fr.wikipedia.org/wiki/Trace_%28alg%C3%A8bre%29})
\end{refs}

\bigskip

\begin{problem}[{Des matrices dont le carré est nul --- Lycée, short}]
\textbf{LIN-23.} \textbf{Problème.} Déterminer, avec démonstration, toutes les matrices réelles
$2 \times 2$ $N$ vérifiant $N^2 = 0$. Montrer qu'une telle matrice $N$ non nulle écrase le
plan tout entier sur une droite passant par l'origine, puis écrase cette droite sur l'origine.
\end{problem}

\noindent Un nombre non nul ne peut avoir un carré nul, mais une matrice non nulle le peut ---
premier signe que les matrices sont une arithmétique véritablement neuve plutôt que des nombres
déguisés. De telles matrices sont dites \textit{nilpotentes} ; le plus petit exemple,
$\begin{pmatrix} 0 & 1 \\ 0 & 0 \end{pmatrix}$, est le cisaillement qui pousse le plan de
côté, et c'est l'unique brique élémentaire dont est assemblée la réduction de Jordan de LIN-14.

\begin{refs}
\item Contexte : Matrice nilpotente --- Wikipédia (\url{https://fr.wikipedia.org/wiki/Matrice_nilpotente})
\item Lecture : G. Strang, \textit{Introduction to Linear Algebra}, ch. 2
\end{refs}

\bigskip

\begin{problem}[{Combien y a-t-il de carrés magiques, vraiment ? --- Lycée}]
\textbf{LIN-24.} \textbf{Problème.} Appelons \textit{carré magique} un tableau $3 \times 3$ de nombres réels dont les
trois sommes par ligne, les trois sommes par colonne et les deux sommes diagonales sont toutes égales
à un même nombre $S$, la somme magique. Les coefficients peuvent être des réels quelconques,
positifs ou non.
\begin{enumerate}
\item[(a)] Démontrer que l'ensemble $M$ de tous les carrés magiques $3 \times 3$ est un espace
vectoriel : une somme de carrés magiques est magique, et tout multiple scalaire aussi.
\item[(b)] Démontrer que le coefficient central d'un carré magique vaut toujours $S/3$.
\item[(c)] Déterminer $\dim M$, avec démonstration, et écrire explicitement une base. (De façon
équivalente : combien de coefficients peut-on choisir librement avant que les autres soient
imposés ?)
\item[(d)] Démontrer que tout carré magique dont les neuf coefficients sont exactement les entiers
$1, 2, \dots, 9$ s'obtient à partir du carré Lo Shu
$\begin{pmatrix} 4 & 9 & 2 \\ 3 & 5 & 7 \\ 8 & 1 & 6 \end{pmatrix}$
par rotation ou réflexion de la grille --- il y en a donc exactement huit.
\end{enumerate}

\end{problem}

\noindent Les carrés magiques sont anciens --- le motif Lo Shu apparaît dans des textes chinois
vieux de plus de deux mille ans, et Dürer en plaça un $4 \times 4$ dans \textit{Melencolia I} (1514),
datant la gravure dans sa dernière ligne. Pendant presque toute cette histoire, on les a traités
comme des curiosités. La lecture par l'algèbre linéaire change entièrement la question : les huit
conditions magiques sont des équations linéaires homogènes, donc les carrés magiques forment le
noyau d'une application linéaire, et demander \og{} combien y en a-t-il \fg{} revient à demander une
dimension. La partie (iv) est le dénombrement classique, et il est frappant que la réponse soit huit
plutôt que huit mille.

\begin{refs}
\item Contexte : Carré magique (mathématiques) --- Wikipédia (\url{https://fr.wikipedia.org/wiki/Carr%C3%A9_magique_%28math%C3%A9matiques%29})
\item Contexte : Noyau (algèbre) --- Wikipédia (\url{https://fr.wikipedia.org/wiki/Noyau_%28alg%C3%A8bre%29})
\item Lecture : G. Strang, \textit{Introduction to Linear Algebra}, ch. 3
\end{refs}

\bigskip

\begin{problem}[{Le produit vectoriel, et le volume qu'il mesure --- Lycée}]
\textbf{LIN-34.} \textbf{Problème.} Pour des vecteurs $u = (u_1, u_2, u_3)$ et $v = (v_1, v_2, v_3)$ de l'espace, on pose
\[ u \times v = (u_2 v_3 - u_3 v_2,\; u_3 v_1 - u_1 v_3,\; u_1 v_2 - u_2 v_1). \]
\begin{enumerate}
\item[(a)] Démontrer que $u \times v$ est orthogonal à la fois à $u$ et à $v$, et démontrer l'identité
de Lagrange
\[ |u \times v|^2 = |u|^2 |v|^2 - (u \cdot v)^2 . \]
En déduire que $|u \times v| = |u|\,|v| \sin\theta$, l'aire du parallélogramme engendré par $u$
et $v$, et que $u \times v = 0$ exactement lorsque $u$ et $v$ sont colinéaires.
\item[(b)] Démontrer que le produit mixte $u \cdot (v \times w)$ est égal au déterminant de la matrice
$3 \times 3$ de lignes $u, v, w$, que sa valeur absolue est le volume du parallélépipède qu'ils
engendrent, et qu'il change de signe lorsqu'on échange deux des trois vecteurs.
\item[(c)] Démontrer que le produit vectoriel n'est pas associatif, en trouvant $u, v, w$ tels que
$(u \times v) \times w \ne u \times (v \times w)$. Démontrer ensuite le développement
$u \times (v \times w) = v\,(u \cdot w) - w\,(u \cdot v)$ et l'identité de Jacobi
\[ u \times (v \times w) + v \times (w \times u) + w \times (u \times v) = 0 . \]
\item[(d)] En déduire que trois vecteurs appartiennent à un même plan passant par l'origine si et seulement
si leur produit mixte est nul, et donner une explication purement géométrique de la raison pour
laquelle le déterminant de la partie (b) devait être un volume.
\end{enumerate}

\end{problem}

\noindent Le produit vectoriel est un fragment d'une structure plus vaste. En 1843, Hamilton
inventa les quaternions et, avec eux, les mots \textit{vecteur} et \textit{scalaire} ; multiplier deux
quaternions \og{} purs \fg{} produit à la fois le produit scalaire et le produit vectoriel, l'un comme partie
scalaire et l'autre comme partie vectorielle. Grassmann disposait dès 1844 d'un produit plus général,
et fut ignoré. Ce sont Gibbs, dans des notes de cours imprimées à compte d'auteur en 1881, et
Heaviside indépendamment, qui scindèrent le produit de Hamilton en les deux opérations familières ---
provoquant une querelle publique d'une décennie avec les quaternionistes --- et c'est E. B. Wilson,
étudiant de Gibbs, qui transforma ces notes en le manuel \textit{Vector Analysis} de 1901, lequel fixa
la notation encore en usage aujourd'hui. Une curiosité qui mérite d'être connue : un produit
bilinéaire de deux vecteurs obéissant aux règles de (a) n'existe qu'en dimensions trois et sept,
celui de dimension sept provenant des octonions.

\begin{refs}
\item Contexte : Produit vectoriel --- Wikipédia (\url{https://fr.wikipedia.org/wiki/Produit_vectoriel})
\item Contexte : Produit mixte --- Wikipédia (\url{https://fr.wikipedia.org/wiki/Produit_mixte})
\item Lecture : E. B. Wilson, \textit{Vector Analysis} (1901), d'après les cours de J. W. Gibbs, ch. 2
\item Lecture : M. J. Crowe, \textit{A History of Vector Analysis}
\end{refs}

\bigskip

\begin{problem}[{Des lignes dont la somme vaut un --- Lycée, short}]
\textbf{LIN-35.} \textbf{Problème.} Soit $P$ une matrice $n \times n$ à coefficients positifs
ou nuls dont chaque ligne a pour somme $1$. Démontrer que le vecteur dont toutes les coordonnées
valent un est vecteur propre de $P$ pour la valeur propre $1$, et démontrer que toute valeur
propre réelle $\lambda$ de $P$ vérifie $|\lambda| \le 1$.
\end{problem}

\noindent Une telle matrice est la table de transition d'une marche aléatoire : le
coefficient $(i,j)$ est la probabilité de sauter de l'état $i$ à l'état $j$. Andreï Markov
introduisit ces chaînes en 1906 pour montrer que la loi des grands nombres survit à la dépendance
entre les épreuves, et il éprouva sa théorie en comptant l'alternance des voyelles et des consonnes
dans \textit{Eugène Onéguine} de Pouchkine. La valeur propre $1$ que vous trouvez ici est la raison
pour laquelle une longue marche aléatoire finit par se stabiliser sur une distribution fixe, et la
version renforcée de cet énoncé --- Perron--Frobenius --- est ce qui classe les pages web en LIN-15.

\begin{refs}
\item Contexte : Matrice stochastique --- Wikipédia (\url{https://fr.wikipedia.org/wiki/Matrice_stochastique})
\item Contexte : Chaîne de Markov --- Wikipédia (\url{https://fr.wikipedia.org/wiki/Cha%C3%AEne_de_Markov})
\end{refs}

\bigskip


\section*{Licence}

\begin{problem}[{Pourquoi la dimension a un sens --- Licence}]
\textbf{LIN-04.} \textbf{Problème.} Des vecteurs $v_1, \dots, v_k$ d'un espace vectoriel sont
\textit{linéairement indépendants} si $c_1 v_1 + \cdots + c_k v_k = 0$ force tous les $c_i = 0$ ;
ils sont \textit{générateurs} si tout vecteur est une telle combinaison ; une \textit{base} est les
deux à la fois.
\begin{enumerate}
\item[(a)] Démontrer le
lemme d'échange de Steinitz : si $v_1, \dots, v_m$ sont indépendants et si $w_1, \dots, w_n$
engendrent, alors $m \le n$.
\item[(b)] En conclure que deux bases quelconques d'un espace vectoriel
de type fini ont le même cardinal --- la \textit{dimension} est donc bien définie.
\item[(c)] Démontrer que toute famille libre
s'étend en une base, et que toute famille génératrice en contient une.
\item[(d)] Montrer que les nombres réels forment
un espace vectoriel de dimension infinie sur $\mathbb{Q}$.
\end{enumerate}

\end{problem}

\noindent On a employé le mot \og{} dimension \fg{} avec assurance pendant un siècle avant que
quiconque démontre qu'il était bien défini ; le lemme d'échange qui règle la question remonte à
l'\textit{Ausdehnungslehre} (1844) de Grassmann, restée négligée, et fut formalisé par Steinitz vers
1910. Le livre de Grassmann --- ignoré de son vivant, son auteur réduit à écrire sur le sanskrit ---
contenait cinquante ans à l'avance l'essentiel de l'algèbre linéaire abstraite. L'espace vectoriel
axiomatique lui-même est dû à Peano (1888).

\begin{refs}
\item Contexte : Indépendance linéaire --- Wikipédia (\url{https://fr.wikipedia.org/wiki/Ind%C3%A9pendance_lin%C3%A9aire})
\item Contexte : Lemme de Steinitz --- Wikipédia (\url{https://fr.wikipedia.org/wiki/Lemme_de_Steinitz})
\item Lecture : S. Axler, \textit{Linear Algebra Done Right}, ch. 2
\end{refs}

\bigskip

\begin{problem}[{Le théorème du rang --- Licence}]
\textbf{LIN-05.} \textbf{Problème.} Soit $T : V \to W$ une application linéaire entre
espaces vectoriels de dimension finie.
\begin{enumerate}
\item[(a)] Démontrer le théorème du rang :
$\dim V = \dim \ker T + \dim \operatorname{im} T$.
\item[(b)] En déduire qu'un système homogène de
$m$ équations linéaires à $n > m$ inconnues admet toujours une solution non nulle.
\item[(c)] En déduire
que, pour une matrice carrée, injectivité, surjectivité et inversibilité sont équivalentes.
\item[(d)] Démontrer que le rang par lignes est égal au rang par colonnes, pour toute matrice.
\end{enumerate}

\end{problem}

\noindent Le théorème du rang est la loi de conservation de l'algèbre linéaire : les
dimensions ne sont jamais créées ni détruites, seulement redistribuées entre noyau et image. Ses
ancêtres sont les arguments de dénombrement d'Euler et la théorie des \og{} diviseurs élémentaires \fg{} de
Frobenius dans les années 1870 ; l'énoncé net en termes d'applications est du vingtième siècle,
popularisé par les traités axiomatiques de l'école de Noether. La partie (iv) --- que personne ne
trouve évidente à première vue --- a été appelée par Strang \og{} le premier grand théorème de l'algèbre
linéaire \fg{}.

\begin{refs}
\item Contexte : Théorème du rang --- Wikipédia (\url{https://fr.wikipedia.org/wiki/Th%C3%A9or%C3%A8me_du_rang})
\item Contexte : Rang (algèbre linéaire) --- Wikipédia (\url{https://fr.wikipedia.org/wiki/Rang_%28alg%C3%A8bre_lin%C3%A9aire%29})
\item Lecture : S. Axler, \textit{Linear Algebra Done Right}, ch. 3
\item Cours : MIT OCW 18.06 Linear Algebra (en anglais) (\url{https://ocw.mit.edu/courses/18-06-linear-algebra-spring-2010/})
\end{refs}

\bigskip

\begin{problem}[{Le déterminant, axiomatiquement --- Licence}]
\textbf{LIN-06.} \textbf{Problème.}
\begin{enumerate}
\item[(a)] Démontrer qu'il existe exactement une fonction $\det$ sur les
matrices réelles $n \times n$ qui soit linéaire en chaque ligne séparément, change de signe
lorsqu'on échange deux lignes, et vérifie $\det I = 1$.
\item[(b)] À partir de ces seuls axiomes, démontrer
$\det(AB) = \det A \cdot \det B$ et $\det A^{\mathsf{T}} = \det A$.
\item[(c)] Démontrer la
signification géométrique : une application linéaire $T$ de $\mathbb{R}^n$ multiplie le volume de
tout pavé par $|\det T|$.
\item[(d)] Expliquer, avec démonstration, pourquoi une matrice est inversible exactement lorsque son
déterminant est non nul.
\end{enumerate}

\end{problem}

\noindent Dès 1812, Cauchy avait organisé les déterminants en une théorie ; vers 1900, ils
dominaient si complètement les manuels que les espaces vectoriels n'y étaient plus qu'une
arrière-pensée. La caractérisation axiomatique de (i), due pour l'essentiel à Weierstrass et
popularisée par les belles notes de cours d'Artin \textit{Galois Theory} ainsi que par Strang, explique
pourquoi une seule et même formule ne cesse de reparaître dans les questions de volume,
d'orientation et d'inversibilité : il n'existe qu'une fonction possible. Le manuel d'Axler bannit
fameusement les déterminants au dernier chapitre pour prouver quelque chose ; ce problème est
l'argument inverse.

\begin{refs}
\item Contexte : Déterminant (mathématiques) --- Wikipédia (\url{https://fr.wikipedia.org/wiki/D%C3%A9terminant_%28math%C3%A9matiques%29})
\item Lecture : M. Artin, \textit{Algebra}, ch. 1
\item Lecture : G. Strang, \textit{Introduction to Linear Algebra}, ch. 5
\item Cours : MIT OCW 18.06 Linear Algebra (en anglais) (\url{https://ocw.mit.edu/courses/18-06-linear-algebra-spring-2010/})
\end{refs}

\bigskip

\begin{problem}[{Valeurs propres, vecteurs propres et Fibonacci --- Licence}]
\textbf{LIN-07.} \textbf{Problème.} Un vecteur non nul $v$ est \textit{vecteur propre} de la
matrice carrée $A$ pour la \textit{valeur propre} $\lambda$ si $Av = \lambda v$.
\begin{enumerate}
\item[(a)] Démontrer que
les valeurs propres de $A$ sont exactement les racines du polynôme caractéristique
$\det(A - \lambda I)$.
\item[(b)] Démontrer que des vecteurs propres associés à des valeurs propres distinctes sont
linéairement indépendants.
\item[(c)] Les nombres de Fibonacci vérifient
$\begin{pmatrix} F_{n+1} \\ F_n \end{pmatrix} = A \begin{pmatrix} F_n \\ F_{n-1} \end{pmatrix}$
pour la matrice $A = \begin{pmatrix} 1 & 1 \\ 1 & 0 \end{pmatrix}$. Diagonaliser $A$
et en déduire une formule close pour $F_n$ en fonction du nombre d'or.
\end{enumerate}

\end{problem}

\noindent Les valeurs propres n'ont pas été trouvées dans les matrices mais dans le ciel :
étudiant les perturbations séculaires des orbites planétaires, Lagrange et Laplace (années 1770--80)
rencontrèrent l'équation caractéristique et eurent besoin que ses racines soient réelles pour la
stabilité du système solaire. Les mots \textit{eigenwert} et \textit{eigenvektor} sont de Hilbert (1904).
La partie (iii) --- la formule de Binet, en réalité connue de de Moivre dès 1730 --- reste la plus nette
des premières démonstrations que la diagonalisation transforme une dynamique en arithmétique.

\begin{refs}
\item Contexte : Valeur propre, vecteur propre et espace propre --- Wikipédia (\url{https://fr.wikipedia.org/wiki/Valeur_propre%2C_vecteur_propre_et_espace_propre})
\item Contexte : Polynôme caractéristique --- Wikipédia (\url{https://fr.wikipedia.org/wiki/Polyn%C3%B4me_caract%C3%A9ristique})
\item Lecture : G. Strang, \textit{Introduction to Linear Algebra}, ch. 6
\item Cours : MIT OCW 18.06 Linear Algebra (en anglais) (\url{https://ocw.mit.edu/courses/18-06-linear-algebra-spring-2010/})
\end{refs}

\bigskip

\begin{problem}[{Quand une matrice est-elle diagonalisable ? --- Licence}]
\textbf{LIN-08.} \textbf{Problème.} Une matrice carrée est \textit{diagonalisable} si elle est semblable
à une matrice diagonale, c'est-à-dire $A = PDP^{-1}$.
\begin{enumerate}
\item[(a)] Démontrer que $A$ (sur $\mathbb{C}$) est
diagonalisable si et seulement si son polynôme minimal --- le polynôme unitaire de plus bas degré
$m$ tel que $m(A) = 0$ --- n'a pas de racine multiple.
\item[(b)] Montrer que
$\begin{pmatrix} 1 & 1 \\ 0 & 1 \end{pmatrix}$ n'est pas diagonalisable, et décrire
géométriquement ce qu'elle fait à la place.
\item[(c)] Démontrer que toute famille commutative de matrices
diagonalisables est \textit{simultanément} diagonalisable.
\item[(d)] Montrer que les matrices diagonalisables sont
denses : toute matrice carrée complexe est limite de matrices diagonalisables.
\end{enumerate}

\end{problem}

\noindent Savoir si une matrice est diagonalisable, c'est toute la question de savoir à quel
point une application linéaire peut être rendue simple, et les cas d'échec de (ii) ne sont pas des
pathologies mais des cisaillements --- ils apparaissent dans tout calcul de déformations. Le critère
(i) descend de la théorie des diviseurs élémentaires de Weierstrass et Frobenius (années 1870). La
partie (iii) a un poids physique : en mécanique quantique, les observables qui commutent sont
exactement celles que l'on peut mesurer simultanément. La partie (iv) est l'astuce préférée de
l'analyste pour démontrer des identités matricielles \og{} d'abord dans le cas générique \fg{}.

\begin{refs}
\item Contexte : Matrice diagonalisable --- Wikipédia (\url{https://fr.wikipedia.org/wiki/Matrice_diagonalisable})
\item Contexte : Polynôme minimal d'un endomorphisme --- Wikipédia (\url{https://fr.wikipedia.org/wiki/Polyn%C3%B4me_minimal_d%27un_endomorphisme})
\item Lecture : S. Axler, \textit{Linear Algebra Done Right}, ch. 5
\end{refs}

\bigskip

\begin{problem}[{Cayley--Hamilton, honnêtement --- Licence}]
\textbf{LIN-09.} \textbf{Problème.} Soit $p(\lambda) = \det(\lambda I - A)$ le
polynôme caractéristique d'une matrice $n \times n$ $A$.
\begin{enumerate}
\item[(a)] Le théorème de Cayley--Hamilton
affirme que $p(A) = 0$. Expliquer précisément pourquoi la tentante \og{} démonstration \fg{} en une ligne
--- substituer $\lambda = A$ dans $\det(\lambda I - A)$ pour obtenir $\det(A - A) = 0$ --- n'a
aucun sens.
\item[(b)] Donner une démonstration correcte, sur $\mathbb{C}$. (Une voie : la démontrer pour les matrices
diagonalisables et utiliser leur densité, LIN-08(iv) ; une autre : l'identité de la comatrice.)
(iii) En déduire que, pour $A$ inversible, l'inverse $A^{-1}$ est un polynôme en $A$ de degré au
plus $n - 1$. (iv) Le théorème vaut-il sur un anneau commutatif quelconque ? Démontrer ou
réfuter.
\end{enumerate}

\end{problem}

\noindent Cayley énonça le théorème dans son mémoire de 1858, le vérifia pour les matrices
$2 \times 2$ et $3 \times 3$, et écrivit --- dans l'un des grands haussements d'épaules des
mathématiques --- qu'il n'avait \og{} pas jugé nécessaire d'entreprendre le labeur d'une démonstration
formelle [\dots{}] dans le cas général \fg{}. Hamilton avait vérifié sa propre version pour les quaternions ;
la première démonstration générale est celle de Frobenius (1878). L'argument fallacieux de (i) a
piégé des générations d'étudiants, ce qui est exactement pourquoi il vaut la peine d'être disséqué :
comprendre pourquoi il échoue apprend ce qu'est vraiment un polynôme d'une matrice.

\begin{refs}
\item Contexte : Théorème de Cayley-Hamilton --- Wikipédia (\url{https://fr.wikipedia.org/wiki/Th%C3%A9or%C3%A8me_de_Cayley-Hamilton})
\item Lecture : S. Axler, \textit{Linear Algebra Done Right}, ch. 8
\item Lecture : R. Horn \& C. Johnson, \textit{Matrix Analysis}, 2e éd., ch. 2
\end{refs}

\bigskip

\begin{problem}[{Le théorème spectral pour les matrices symétriques --- Licence}]
\textbf{LIN-10.} \textbf{Problème.} Soit $A$ une matrice réelle symétrique $n \times n$
($A = A^{\mathsf{T}}$). Démontrer :
\begin{enumerate}
\item[(a)] toute valeur propre de $A$ est réelle ;
\item[(b)] des vecteurs propres associés à des
valeurs propres distinctes sont orthogonaux ;
\item[(c)] $\mathbb{R}^n$ admet une base orthonormée de
vecteurs propres de $A$, c'est-à-dire $A = Q \Lambda Q^{\mathsf{T}}$ avec $Q$ orthogonale et
$\Lambda$ diagonale.
\item[(d)] En déduire le théorème des axes principaux : toute forme quadratique
$q(x) = x^{\mathsf{T}} A x$ devient une somme pondérée de carrés
$\lambda_1 y_1^2 + \cdots + \lambda_n y_n^2$ dans des coordonnées convenablement tournées, et toute
ellipse ou hyperbole $q(x) = 1$ a des axes perpendiculaires.
\end{enumerate}

\end{problem}

\noindent Cauchy démontra en 1829 le théorème sur la réalité des valeurs propres, motivé
exactement par (iv) : les axes principaux des quadriques et la stabilité des corps en rotation. C'est
le théorème le plus utilisé du domaine --- la statistique (analyse en composantes principales), la
mécanique (modes propres), la théorie quantique (observables) reposent toutes dessus --- et son
extension en dimension infinie, portée par Hilbert et von Neumann, a créé l'analyse fonctionnelle.
Toute démonstration doit quelque part utiliser deux fois la symétrie ; trouver où fait partie de
l'exercice.

\begin{refs}
\item Contexte : Théorème spectral --- Wikipédia (\url{https://fr.wikipedia.org/wiki/Th%C3%A9or%C3%A8me_spectral})
\item Lecture : S. Axler, \textit{Linear Algebra Done Right}, ch. 7
\item Lecture : G. Strang, \textit{Introduction to Linear Algebra}, ch. 6
\item Cours : MIT OCW 18.06 Linear Algebra (en anglais) (\url{https://ocw.mit.edu/courses/18-06-linear-algebra-spring-2010/})
\end{refs}

\bigskip

\begin{problem}[{Quatre visages de la définie-positivité --- Licence}]
\textbf{LIN-11.} \textbf{Problème.} Une matrice réelle symétrique $A$ est \textit{définie positive}
si $x^{\mathsf{T}} A x > 0$ pour tout vecteur non nul $x$. Démontrer que les assertions
suivantes sont équivalentes :
\begin{enumerate}
\item[(a)] $A$ est définie positive ;
\item[(b)] toutes les valeurs propres de $A$ sont strictement positives ;
\item[(c)] critère de Sylvester : les $n$ mineurs principaux dominants (déterminants des blocs
$k \times k$ en haut à gauche) sont tous strictement positifs ;
\item[(d)] $A = B^{\mathsf{T}} B$ pour une matrice inversible
$B$ --- de façon équivalente, $A$ admet une factorisation de Cholesky $A = LL^{\mathsf{T}}$ avec
$L$ triangulaire inférieure à diagonale strictement positive. Exhiber ensuite une matrice symétrique
dont tous les coefficients diagonaux sont strictement positifs et qui n'est \textit{pas} définie
positive.
\end{enumerate}

\end{problem}

\noindent La définie-positivité est le lieu où l'algèbre linéaire rencontre la notion
physique d'énergie : $x^{\mathsf{T}} A x$ est l'énergie emmagasinée dans une configuration, et (d)
dit que toute énergie de ce type est une somme de carrés. Le critère de Sylvester date de la théorie
des formes quadratiques des années 1850 ; la factorisation de Cholesky fut conçue par un officier
d'artillerie français pour des levés géodésiques et ne fut publiée qu'après sa mort pendant la
Première Guerre mondiale. Tester la définie-positivité est aujourd'hui une primitive de base de
l'optimisation --- la définie-positivité d'une hessienne, c'est ce que signifie le \og{} test de la dérivée
seconde \fg{} en dimension $n$.

\begin{refs}
\item Contexte : Matrice définie positive --- Wikipédia (\url{https://fr.wikipedia.org/wiki/Matrice_d%C3%A9finie_positive})
\item Contexte : Critère de Sylvester --- Wikipédia (\url{https://fr.wikipedia.org/wiki/Matrice_d%C3%A9finie_positive#Crit%C3%A8re_de_Sylvester})
\item Contexte : Factorisation de Cholesky --- Wikipédia (\url{https://fr.wikipedia.org/wiki/Factorisation_de_Cholesky})
\item Lecture : R. Horn \& C. Johnson, \textit{Matrix Analysis}, 2e éd., ch. 7
\end{refs}

\bigskip

\begin{problem}[{L'exponentielle d'une matrice --- Licence}]
\textbf{LIN-12.} \textbf{Problème.} Pour une matrice carrée $A$, on pose
$e^{A} = \sum_{k=0}^{\infty} A^k / k!$.
\begin{enumerate}
\item[(a)] Démontrer que la série converge pour toute $A$.
\item[(b)] Démontrer que si $AB = BA$ alors $e^{A+B} = e^A e^B$, et trouver deux matrices pour lesquelles
$e^{A+B} \ne e^A e^B$.
\item[(c)] Démontrer $\det e^A = e^{\operatorname{tr} A}$.
\item[(d)] Montrer que
$x(t) = e^{tA} x_0$ est l'unique solution de l'équation différentielle $x' = Ax$ avec
$x(0) = x_0$, et calculer explicitement $e^{tA}$ lorsque $A = \begin{pmatrix} 0 & -1 \\
1 & 0 \end{pmatrix}$. Quelles fonctions familières apparaissent, et pourquoi le devaient-elles ?
\end{enumerate}

\end{problem}

\noindent L'exponentielle de matrice est la façon dont l'algèbre linéaire parle à l'analyse :
elle convertit une matrice statique en un flot, et elle sous-tend toute équation différentielle à
coefficients constants, tout semi-groupe de Markov, et l'ensemble de la théorie des groupes de Lie
(où l'échec de (ii) pour des matrices non commutantes n'est pas un défaut mais le début de
l'histoire). La calculer de façon fiable est notoirement délicat --- l'article de synthèse de Moler et
Van Loan, \og{} Nineteen Dubious Ways to Compute the Exponential of a Matrix \fg{}, est un classique
précisément parce que les dix-neuf méthodes peuvent toutes échouer.

\begin{refs}
\item Contexte : Exponentielle d'une matrice --- Wikipédia (\url{https://fr.wikipedia.org/wiki/Exponentielle_d%27une_matrice})
\item Lecture : G. Strang, \textit{Introduction to Linear Algebra}, ch. 6
\item Article : C. Moler \& C. Van Loan, \og{} Nineteen Dubious Ways to Compute the Exponential of a Matrix, Twenty-Five Years Later \fg{} (2003), SIAM Review 45, 3--49
\end{refs}

\bigskip

\begin{problem}[{La décomposition en valeurs singulières --- Licence}]
\textbf{LIN-13.} \textbf{Problème.} (i) Démontrer que toute matrice réelle $m \times n$ $A$ peut
s'écrire $A = U \Sigma V^{\mathsf{T}}$, où $U$ et $V$ sont orthogonales et $\Sigma$
est diagonale à coefficients positifs ou nuls $\sigma_1 \ge \sigma_2 \ge \cdots \ge 0$ (les
\textit{valeurs singulières}). (Indication pour l'existence : appliquer le théorème spectral, LIN-10, à
$A^{\mathsf{T}} A$.) (ii) Montrer que, géométriquement, $A$ envoie la sphère unité sur un
ellipsoïde dont les demi-axes ont pour longueurs les $\sigma_i$. (iii) Démontrer le théorème
d'Eckart--Young : parmi toutes les matrices de rang au plus $k$, la troncature
$A_k = \sum_{i \le k} \sigma_i u_i v_i^{\mathsf{T}}$ est la plus proche de $A$ pour la norme
spectrale, avec une erreur $\sigma_{k+1}$.
\end{problem}

\noindent La décomposition en valeurs singulières fut découverte par Beltrami (1873) et
Camille Jordan (1874) comme forme canonique des formes bilinéaires, puis refondée comme outil de
calcul par Golub et Kahan en 1965. La partie (iii), démontrée par Eckart et Young en 1936 dans une
revue de psychométrie, est sans doute le théorème le plus rentable du vingtième siècle :
l'approximation optimale de rang faible est le moteur mathématique de la compression d'images, des
systèmes de recommandation, du filtrage du bruit et de l'analyse en composantes principales.
Trefethen et Bau ouvrent leur ouvrage de référence par la décomposition en valeurs singulières, et
non par le déterminant.

\begin{refs}
\item Contexte : Décomposition en valeurs singulières --- Wikipédia (\url{https://fr.wikipedia.org/wiki/D%C3%A9composition_en_valeurs_singuli%C3%A8res})
\item Lecture : L. Trefethen \& D. Bau, \textit{Numerical Linear Algebra}, leçons 4--5
\item Article : C. Eckart \& G. Young, \og{} The approximation of one matrix by another of lower rank \fg{} (1936), Psychometrika 1, 211--218
\item Cours : MIT OCW 18.065 Matrix Methods in Data Analysis (en anglais) (\url{https://ocw.mit.edu/courses/18-065-matrix-methods-in-data-analysis-signal-processing-and-machine-learning-spring-2018/})
\end{refs}

\bigskip

\begin{problem}[{Gram--Schmidt est une factorisation --- Licence, short}]
\textbf{LIN-25.} \textbf{Problème.} Soit $A$ une matrice réelle $m \times n$ avec $m \ge n$
dont les colonnes sont linéairement indépendantes. Démontrer que $A = QR$ où $Q$ est
$m \times n$ à colonnes orthonormées et $R$ est $n \times n$ triangulaire supérieure à
coefficients diagonaux strictement positifs, et démontrer que $Q$ et $R$ sont uniquement
déterminées par $A$.
\end{problem}

\noindent Gram--Schmidt s'enseigne d'ordinaire comme une procédure --- orthogonaliser chaque
colonne contre les précédentes --- et l'intérêt de ce problème est que la comptabilité de cette
procédure \textit{est} une factorisation matricielle, $R$ enregistrant chaque coefficient utilisé. Le
procédé porte le nom de Jørgen Gram (1883) et d'Erhard Schmidt (1907), bien que Laplace l'eût employé
presque un siècle plus tôt. En pratique numérique, l'algorithme classique perd gravement
l'orthogonalité, et l'on calcule $QR$ à la place par réflexions de Householder ; la factorisation
obtenue sous-tend à la fois les moindres carrés (LIN-27) et l'algorithme QR pour les valeurs
propres.

\begin{refs}
\item Contexte : Décomposition QR --- Wikipédia (\url{https://fr.wikipedia.org/wiki/D%C3%A9composition_QR})
\item Contexte : Algorithme de Gram-Schmidt --- Wikipédia (\url{https://fr.wikipedia.org/wiki/Algorithme_de_Gram-Schmidt})
\item Lecture : L. Trefethen \& D. Bau, \textit{Numerical Linear Algebra}, leçons 7--10
\end{refs}

\bigskip

\begin{problem}[{Le quotient de Rayleigh est pris au piège --- Licence, short}]
\textbf{LIN-26.} \textbf{Problème.} Soit $A$ une matrice réelle symétrique $n \times n$ de plus
grande valeur propre $\lambda_{\max}$ et de plus petite valeur propre $\lambda_{\min}$. Démontrer
que, pour tout $x \in \mathbb{R}^n$ non nul,
\[ \lambda_{\min} \;\le\; \frac{x^{\mathsf{T}} A x}{x^{\mathsf{T}} x} \;\le\; \lambda_{\max}, \]
et que l'égalité de part ou d'autre a lieu exactement lorsque $x$ est vecteur propre pour la valeur
propre correspondante.
\end{problem}

\noindent Lord Rayleigh introduisit ce quotient dans \textit{The Theory of Sound} (1877)
pour estimer les fréquences d'un corps vibrant : devinez une forme, calculez le quotient, et vous
tenez une borne supérieure de la note la plus grave. Le problème convertit une question algébrique ---
les racines d'un polynôme --- en une optimisation, ce qui explique pourquoi les valeurs propres des
matrices symétriques sont stables et calculables ; l'énoncé général pour la $k$-ième valeur propre
est le théorème de Courant--Fischer de LIN-16.

\begin{refs}
\item Contexte : Quotient de Rayleigh --- Wikipédia (\url{https://fr.wikipedia.org/wiki/Quotient_de_Rayleigh})
\item Lecture : R. Horn \& C. Johnson, \textit{Matrix Analysis}, 2e éd., ch. 4
\end{refs}

\bigskip

\begin{problem}[{Les moindres carrés sont une ombre --- Licence}]
\textbf{LIN-27.} \textbf{Problème.} Soient $A$ une matrice réelle $m \times n$ et $b \in \mathbb{R}^m$. Le
problème des moindres carrés consiste à minimiser $\|Ax - b\|$ sur tous les $x \in \mathbb{R}^n$,
où $\|\cdot\|$ est la norme euclidienne.
\begin{enumerate}
\item[(a)] Démontrer que $x$ minimise $\|Ax - b\|$ si et seulement si $Ax$ est la projection orthogonale
de $b$ sur l'espace des colonnes de $A$, et que cela a lieu si et seulement si $x$ résout les
équations normales
\[ A^{\mathsf{T}} A x = A^{\mathsf{T}} b. \]
En conclure qu'un minimiseur existe toujours, quels que soient $A$ et $b$.
\item[(b)] Démontrer que le minimiseur est unique exactement lorsque les colonnes de $A$ sont indépendantes,
et que dans ce cas la projection sur l'espace des colonnes est
$P = A (A^{\mathsf{T}} A)^{-1} A^{\mathsf{T}}$. Vérifier $P^{\mathsf{T}} = P$ et $P^2 = P$,
et montrer que $I - P$ est la projection orthogonale sur $\ker A^{\mathsf{T}}$.
\item[(c)] Démontrer la réciproque de (ii) : toute matrice symétrique idempotente est la projection orthogonale
sur son propre espace des colonnes.
\item[(d)] À l'aide de la factorisation $A = QR$ de LIN-25, montrer que les équations normales deviennent
$Rx = Q^{\mathsf{T}} b$, un système triangulaire. Expliquer pourquoi former $A^{\mathsf{T}} A$
élève au carré le conditionnement de $A$, et donc pourquoi c'est cette voie-ci que les analystes
numériciens empruntent réellement.
\end{enumerate}

\end{problem}

\noindent Legendre publia la méthode des moindres carrés en 1805, dans un appendice sur la
détermination des orbites cométaires, et lui donna son nom ; Gauss répliqua qu'il l'employait depuis
1795 et qu'il s'en était servi en 1801 pour prédire où reparaîtrait l'astéroïde perdu Cérès --- ce
qu'il fit. La querelle de priorité qui s'ensuivit fut âpre, et les mathématiques lui ont survécu : le
théorème de Gauss--Markov fait des moindres carrés le meilleur estimateur linéaire sans biais, et
aujourd'hui les mêmes équations normales siègent au cœur de la régression, du positionnement GPS et
de l'ajustement de presque tout modèle empirique. Le contenu géométrique tient en une phrase --- la
meilleure approximation est une ombre perpendiculaire --- et tout le reste est comptabilité.

\begin{refs}
\item Contexte : Méthode des moindres carrés --- Wikipédia (\url{https://fr.wikipedia.org/wiki/M%C3%A9thode_des_moindres_carr%C3%A9s})
\item Contexte : Projecteur (mathématiques) --- Wikipédia (\url{https://fr.wikipedia.org/wiki/Projecteur_%28math%C3%A9matiques%29})
\item Lecture : L. Trefethen \& D. Bau, \textit{Numerical Linear Algebra}, leçons 11 et 18--19
\item Cours : MIT OCW 18.06 Linear Algebra (en anglais) (\url{https://ocw.mit.edu/courses/18-06-linear-algebra-spring-2010/})
\end{refs}

\bigskip

\begin{problem}[{Quand a-t-on $A = LU$ ? --- Licence}]
\textbf{LIN-36.} \textbf{Problème.} Exécuter le pivot de Gauss sur une matrice carrée $A$ sans échange de lignes,
en enregistrant les multiplicateurs, revient à écrire $A = LU$ avec $L$ triangulaire inférieure
unitaire (des uns sur la diagonale) et $U$ triangulaire supérieure.
\begin{enumerate}
\item[(a)] Démontrer qu'une matrice inversible $A$ admet une telle factorisation si et seulement si ses
$n$ mineurs principaux dominants --- les déterminants des blocs $k \times k$ en haut à gauche ---
sont tous non nuls, et démontrer que $L$ et $U$ sont alors uniques.
\item[(b)] Exhiber une matrice inversible sans factorisation $LU$. Démontrer ensuite que pour toute matrice
carrée $A$ il existe une matrice de permutation $P$ telle que $PA = LU$ : les échanges de lignes
sauvent toujours l'algorithme.
\item[(c)] Démontrer qu'une matrice symétrique définie positive (LIN-11) n'a jamais besoin d'échange de
lignes, et qu'elle admet une unique factorisation $A = L L^{\mathsf{T}}$ avec $L$ triangulaire
inférieure à diagonale strictement positive --- la factorisation de Cholesky, qui coûte environ la
moitié des opérations du cas général.
\item[(d)] On définit le facteur de croissance de l'élimination comme le plus grand coefficient en valeur
absolue apparaissant à un moment quelconque du processus, divisé par le plus grand coefficient en
valeur absolue de $A$. Démontrer que même avec pivotage partiel --- en choisissant toujours le plus
grand pivot disponible dans la colonne courante --- le facteur de croissance pour une matrice
$n \times n$ peut atteindre $2^{\,n-1}$, en construisant une matrice qui l'atteint. Expliquer
pourquoi cela ne rend l'algorithme standard que conditionnellement stable, et pourquoi il est
néanmoins utilisé partout.
\end{enumerate}

\end{problem}

\noindent L'élimination est ancienne --- le chapitre 8 des \textit{Neuf Chapitres} la pratique ---
mais y voir une \textit{factorisation} est moderne : l'astronome polonais Tadeusz Banachiewicz écrivit
l'élimination comme un produit de matrices triangulaires en 1938, et Alan Turing, dans son article de
1948 \og{} Rounding-off errors in matrix processes \fg{}, donna la notation $LU$ et introduisit le
conditionnement pour expliquer quand on pouvait faire confiance au procédé. L'analyse d'erreur
inverse de Wilkinson dans les années 1960 acheva le tableau. La partie (d) est l'embarras le plus
honnête du domaine : le pire cas est exponentiel, quasiment personne ne l'a jamais rencontré en
pratique, et aucune explication satisfaisante de cet écart n'est connue. Notez aussi à quoi sert la
factorisation : une fois $A = LU$ en main, chaque nouveau second membre coûte $O(n^2)$ plutôt que
$O(n^3)$, et c'est pourquoi personne ne calcule $A^{-1}$.

\begin{refs}
\item Contexte : Décomposition LU --- Wikipédia (\url{https://fr.wikipedia.org/wiki/D%C3%A9composition_LU})
\item Contexte : Factorisation de Cholesky --- Wikipédia (\url{https://fr.wikipedia.org/wiki/Factorisation_de_Cholesky})
\item Lecture : L. Trefethen \& D. Bau, \textit{Numerical Linear Algebra}, leçons 20--23
\item Cours : MIT OCW 18.06 Linear Algebra (en anglais) (\url{https://ocw.mit.edu/courses/18-06-linear-algebra-spring-2010/})
\end{refs}

\bigskip

\begin{problem}[{La pseudo-inverse et la plus petite solution --- Licence}]
\textbf{LIN-37.} \textbf{Problème.} Soit $A$ une matrice réelle $m \times n$. On dit qu'une matrice
$n \times m$ $X$ est une \textit{inverse de Moore--Penrose} de $A$ si
\[ AXA = A, \qquad XAX = X, \qquad (AX)^{\mathsf{T}} = AX, \qquad (XA)^{\mathsf{T}} = XA . \]
\begin{enumerate}
\item[(a)] Démontrer qu'une telle $X$ existe pour toute $A$ et qu'elle est unique ; on la note $A^{+}$.
\item[(b)] Démontrer que si $A = U \Sigma V^{\mathsf{T}}$ est une décomposition en valeurs singulières
(LIN-13), alors $A^{+} = V \Sigma^{+} U^{\mathsf{T}}$, où $\Sigma^{+}$ s'obtient à partir de
$\Sigma$ en inversant les valeurs singulières non nulles et en transposant. En déduire que
$A^{+} = (A^{\mathsf{T}}A)^{-1}A^{\mathsf{T}}$ lorsque les colonnes de $A$ sont indépendantes, et
$A^{+} = A^{-1}$ lorsque $A$ est inversible.
\item[(c)] Démontrer que parmi tous les vecteurs minimisant $\|Ax - b\|$ (LIN-27), le vecteur $x = A^{+}b$
est l'unique de plus petite norme euclidienne --- la pseudo-inverse résout donc le problème des moindres
carrés même lorsque celui-ci n'a pas de solution unique.
\item[(d)] Démontrer que $A \mapsto A^{+}$ n'est pas continue, en exhibant des matrices
$A_\varepsilon \to A$ avec $\|A_\varepsilon^{+}\|$ non bornée. Montrer aussi que
$(AB)^{+} = B^{+}A^{+}$ peut être mise en défaut. Expliquer ce que ces deux faits signifient pour
quiconque calcule une pseudo-inverse en virgule flottante.
\end{enumerate}

\end{problem}

\noindent E. H. Moore définit une \og{} réciproque générale \fg{} de toute matrice en 1920, dans un
résumé destiné à l'American Mathematical Society rédigé dans une notation personnelle que presque
personne ne lut ; Arne Bjerhammar la redécouvrit vers 1951 en géodésie ; et en 1955 Roger Penrose ---
une décennie avant ses travaux sur les trous noirs --- donna la caractérisation par les quatre
conditions ci-dessus, qui reste la définition employée. La pseudo-inverse est ce que les logiciels
statistiques et d'ingénierie calculent silencieusement lorsqu'un système est de rang déficient. La
partie (d) est la raison pour laquelle ces logiciels exigent une tolérance : le rang n'est pas une
fonction continue des coefficients, de sorte qu'une matrice presque de rang déficient a une
pseudo-inverse énorme, et la \og{} solution de norme minimale \fg{} peut sauter de façon discontinue lorsque
les données bougent.

\begin{refs}
\item Contexte : Pseudo-inverse --- Wikipédia (\url{https://fr.wikipedia.org/wiki/Pseudo-inverse})
\item Article : R. Penrose, \og{} A generalized inverse for matrices \fg{} (1955), Proceedings of the Cambridge Philosophical Society 51, 406--413
\item Lecture : G. Golub \& C. Van Loan, \textit{Matrix Computations}, ch. 5
\item Cours : MIT OCW 18.065 Matrix Methods in Data Analysis (en anglais) (\url{https://ocw.mit.edu/courses/18-065-matrix-methods-in-data-analysis-signal-processing-and-machine-learning-spring-2018/})
\end{refs}

\bigskip

\begin{problem}[{Une matrice construite sur mesure --- Licence, short}]
\textbf{LIN-38.} \textbf{Problème.} Pour un polynôme unitaire
$p(x) = x^n + c_{n-1}x^{n-1} + \cdots + c_1 x + c_0$ sur un corps, soit $C(p)$ sa matrice
compagnon : des uns sur la sous-diagonale, les coefficients $-c_0, -c_1, \dots, -c_{n-1}$ le long de
la dernière colonne, et des zéros ailleurs. Démontrer que le polynôme caractéristique et le polynôme
minimal de $C(p)$ sont tous deux égaux à $p$, et en déduire qu'une matrice carrée est semblable à
une matrice compagnon si et seulement si ses polynômes minimal et caractéristique coïncident.
\end{problem}

\noindent Le corollaire immédiat est que \textit{tout} polynôme unitaire est le polynôme
caractéristique de quelqu'un, si bien que le problème des valeurs propres et celui de la recherche
des racines ne sont qu'un seul problème sous deux chapeaux. Cela a une conséquence tranchante :
puisqu'il n'existe pas de formule par radicaux pour les racines d'une quintique générale (ALG-21), il
ne peut exister d'algorithme exact fini pour les valeurs propres d'une matrice $5 \times 5$
générale, et toute routine de calcul de valeurs propres jamais écrite est donc itérative. Le trafic
circule dans les deux sens --- les bibliothèques numériques calculent les racines d'un polynôme comme
les valeurs propres de la matrice compagnon --- et empiler des blocs compagnons donne la forme
canonique rationnelle de Frobenius, cousine indépendante du corps de la réduction de Jordan de
LIN-14.

\begin{refs}
\item Contexte : Matrice compagnon --- Wikipédia (\url{https://fr.wikipedia.org/wiki/Matrice_compagnon})
\item Contexte : Décomposition de Frobenius --- Wikipédia (\url{https://fr.wikipedia.org/wiki/D%C3%A9composition_de_Frobenius})
\item Lecture : R. Horn \& C. Johnson, \textit{Matrix Analysis}, 2e éd., ch. 3
\end{refs}

\bigskip


\section*{Master}

\begin{problem}[{Pourquoi la réduction de Jordan existe --- Master}]
\textbf{LIN-14.} \textbf{Problème.} Soit $T$ un endomorphisme d'un espace vectoriel
complexe de dimension finie $V$.
\begin{enumerate}
\item[(a)] Démontrer que $V$ est la somme directe des \textit{sous-espaces
caractéristiques} $\ker (T - \lambda I)^n$, un pour chaque valeur propre $\lambda$.
\item[(b)] Démontrer
que tout endomorphisme nilpotent $N$ (une de ses puissances est nulle) admet une base dans laquelle
sa matrice est une somme directe de \og{} blocs de décalage \fg{} --- des zéros partout sauf des uns juste
au-dessus de la diagonale --- et que le multiensemble des tailles de blocs est uniquement déterminé par
$N$.
\item[(c)] Combiner (i) et (ii) pour
démontrer le théorème de la réduction de Jordan : $T$ admet une base dans laquelle sa matrice est
une somme directe de blocs de Jordan $\lambda I + \text{décalage}$, unique à l'ordre près.
\item[(d)] En conclure que deux matrices
complexes sont semblables si et seulement si elles ont la même réduite de Jordan, et exhiber deux
matrices ayant les mêmes polynômes caractéristique \textit{et} minimal qui ne sont pas semblables.
\end{enumerate}

\end{problem}

\noindent Camille Jordan démontra la forme canonique dans son \textit{Traité des
substitutions} de 1870 --- pour des matrices sur les corps finis, chose remarquable, le cas complexe
suivant --- en s'appuyant sur la théorie des diviseurs élémentaires de Weierstrass (1868). C'est la
réponse définitive à la question \og{} à quoi ressemblent les applications linéaires ? \fg{} : tout
endomorphisme est diagonal-plus-nilpotent, et la partie nilpotente relève de la combinatoire (une
partition pour chaque valeur propre). La partie (iv) explique pourquoi le théorème est réellement
nécessaire : aucune liste d'invariants polynomiaux en deçà des données complètes de blocs ne classifie
la similitude.

\begin{refs}
\item Contexte : Réduction de Jordan --- Wikipédia (\url{https://fr.wikipedia.org/wiki/R%C3%A9duction_de_Jordan})
\item Contexte : Sous-espace caractéristique --- Wikipédia (\url{https://fr.wikipedia.org/wiki/Sous-espace_caract%C3%A9ristique})
\item Lecture : S. Axler, \textit{Linear Algebra Done Right}, ch. 8
\item Lecture : R. Horn \& C. Johnson, \textit{Matrix Analysis}, 2e éd., ch. 3
\end{refs}

\bigskip

\begin{problem}[{Perron--Frobenius et le vecteur propre à 25 000 000 000 de dollars --- Master}]
\textbf{LIN-15.} \textbf{Problème.}
\begin{enumerate}
\item[(a)] Démontrer le théorème de Perron : une matrice carrée à
coefficients strictement positifs possède une valeur propre réelle $r > 0$ qui est simple,
strictement plus grande en valeur absolue que toute autre valeur propre, et qui possède un vecteur
propre à coefficients strictement positifs --- et aucune autre valeur propre n'a de vecteur propre à
coefficients positifs ou nuls.
\item[(b)] Expliquer par un exemple ce qui peut échouer
lorsque les coefficients sont seulement positifs ou nuls, et énoncer le raffinement de Frobenius pour
les matrices irréductibles.
\item[(c)] PageRank modélise un internaute qui, depuis une page, suit un lien sortant au hasard avec
probabilité $\alpha$ (disons 0,85) et saute sinon vers une page uniformément aléatoire ; le vecteur
de classement est la distribution stationnaire. Formuler la \og{} matrice de Google \fg{} $G$, démontrer via
(i) que le classement existe et est unique, et démontrer que la méthode de la puissance itérée y
converge à une vitesse géométrique gouvernée par $\alpha$.
\end{enumerate}

\end{problem}

\noindent Perron démontra son théorème en 1907, Frobenius l'étendit aux matrices positives
irréductibles en 1912, et pendant quatre-vingt-dix ans il fonctionna discrètement à l'intérieur des
modèles de population, de l'économie entrées-sorties (le prix Nobel de Leontief) et de la théorie des
chaînes de Markov. Puis Brin et Page bâtirent dessus un moteur de recherche. Le titre de l'essai de
Bryan et Leise dans SIAM Review --- \og{} The \$25,000,000,000 eigenvector \fg{} --- chiffrait le théorème à
l'introduction en bourse de Google en 2004 ; l'estimation n'a vieilli que dans un seul sens.

\begin{refs}
\item Contexte : Théorème de Perron-Frobenius --- Wikipédia (\url{https://fr.wikipedia.org/wiki/Th%C3%A9or%C3%A8me_de_Perron-Frobenius})
\item Contexte : PageRank --- Wikipédia (\url{https://fr.wikipedia.org/wiki/PageRank})
\item Article : S. Brin \& L. Page, \og{} The Anatomy of a Large-Scale Hypertextual Web Search Engine \fg{} (1998), Computer Networks and ISDN Systems 30, 107--117
\item Article : K. Bryan \& T. Leise, \og{} The \$25,000,000,000 Eigenvector: The Linear Algebra behind Google \fg{} (2006), SIAM Review 48, 569--581
\end{refs}

\bigskip

\begin{problem}[{Courant--Fischer et l'entrelacement des valeurs propres --- Master}]
\textbf{LIN-16.} \textbf{Problème.} Soit $A$ une matrice réelle symétrique $n \times n$ de
valeurs propres $\lambda_1 \ge \lambda_2 \ge \cdots \ge \lambda_n$. (i) Démontrer le théorème
min--max de Courant--Fischer :
\[ \lambda_k = \max_{\dim U = k} \; \min_{0 \ne x \in U} \; \frac{x^{\mathsf{T}} A x}{x^{\mathsf{T}} x}. \]
(ii) En déduire le théorème d'entrelacement de Cauchy : si $B$ s'obtient à partir de $A$ en
supprimant une ligne et la colonne correspondante, ses valeurs propres
$\mu_1 \ge \cdots \ge \mu_{n-1}$ vérifient
$\lambda_1 \ge \mu_1 \ge \lambda_2 \ge \mu_2 \ge \cdots \ge \mu_{n-1} \ge \lambda_n$.
(iii) En déduire l'inégalité de perturbation de Weyl : pour $A$ et $E$ symétriques, les
$k$-ièmes valeurs propres de $A$ et de $A + E$ diffèrent d'au plus la plus grande valeur propre
en valeur absolue de $E$. Expliquer pourquoi cela fait des valeurs propres des matrices symétriques
des objets \og{} stables \fg{}, en net contraste avec les valeurs propres des matrices générales --- et trouver
un exemple frappant de ce contraste.
\end{problem}

\noindent L'image variationnelle des valeurs propres --- comme solutions de problèmes
d'optimisation plutôt que comme racines de polynômes --- descend des travaux de Rayleigh sur les
vibrations (années 1870) via Fischer (1905) et Courant (1920), et c'est sous cette forme que la
théorie spectrale fait un travail réel : l'entrelacement alimente tout, des suites de Sturm à la
théorie spectrale des graphes, et l'inégalité de Weyl est la raison pour laquelle des fréquences
mesurées signifient quelque chose malgré l'imperfection des appareils. En 2013, les familles
entrelaçantes de polynômes --- une descendante de (ii) --- ont résolu le problème de Kadison--Singer,
vieux de plusieurs décennies (Marcus, Spielman, Srivastava).

\begin{refs}
\item Contexte : Théorème min-max de Courant-Fischer --- Wikipédia (\url{https://fr.wikipedia.org/wiki/Th%C3%A9or%C3%A8me_min-max_de_Courant-Fischer})
\item Lecture : R. Horn \& C. Johnson, \textit{Matrix Analysis}, 2e éd., ch. 4
\item Lecture : L. Trefethen \& D. Bau, \textit{Numerical Linear Algebra}, partie V
\end{refs}

\bigskip

\begin{problem}[{Strassen : multiplier les matrices plus vite qu'on ne vous l'a appris --- Master}]
\textbf{LIN-17.} \textbf{Problème.}
\begin{enumerate}
\item[(a)] Montrer que des matrices $2 \times 2$ sur un anneau commutatif
quelconque peuvent être multipliées en n'utilisant que 7 multiplications d'éléments de l'anneau (au
lieu des 8 évidentes), au prix d'additions supplémentaires.
\item[(b)] Démontrer qu'appliquer cela récursivement à des matrices
par blocs multiplie deux matrices $n \times n$ en $O(n^{\log_2 7}) = O(n^{2{,}81})$ opérations
élémentaires.
\item[(c)] Démontrer que 7 est optimal pour le cas $2 \times 2$ : aucun algorithme bilinéaire n'emploie 6
multiplications (Winograd, 1971 --- difficile ; a minima, lire une démonstration et en reconstruire le
plan).
\item[(d)] Montrer que tout algorithme demande au moins $cn^2$ étapes, de sorte que l'exposant de la
multiplication matricielle est compris entre 2 et $\log_2 7$.
\end{enumerate}

\end{problem}

\noindent Strassen trouva son algorithme en 1969 alors qu'il tentait de démontrer le
contraire --- que le pivot de Gauss est optimal ; son article s'intitule \og{} Gaussian elimination is not
optimal \fg{}. La découverte que l'algorithme cubique \og{} évident \fg{} est battable a ouvert la théorie de la
complexité algébrique et une course de cinquante ans, documentée en LIN-19. La méthode de Strassen
n'est pas une curiosité : pour de grandes matrices, elle est implantée dans des bibliothèques de
production, et sa structure récursive par blocs est le modèle de pratiquement toutes les
multiplications matricielles rapides depuis lors.

\begin{refs}
\item Contexte : Algorithme de Strassen --- Wikipédia (\url{https://fr.wikipedia.org/wiki/Algorithme_de_Strassen})
\item Article : V. Strassen, \og{} Gaussian elimination is not optimal \fg{} (1969), Numerische Mathematik 13, 354--356
\item Lecture : J. Landsberg, \textit{Tensors: Geometry and Applications} (AMS GSM 128), ch. 1
\end{refs}

\bigskip

\begin{problem}[{La porte des tenseurs --- Master}]
\textbf{LIN-18.} \textbf{Problème.}
\begin{enumerate}
\item[(a)] Définir le produit tensoriel $V \otimes W$ de deux espaces
vectoriels par sa propriété universelle : les applications bilinéaires issues de $V \times W$
correspondent exactement aux applications linéaires issues de $V \otimes W$. Démontrer qu'il existe,
qu'il est unique à isomorphisme canonique près, et qu'il vérifie
$\dim(V \otimes W) = \dim V \cdot \dim W$.
\item[(b)] Démontrer que tout élément de $V \otimes W$ n'est pas un tenseur \og{} pur \fg{} $v \otimes w$ ; le
nombre minimal de tenseurs purs nécessaires pour écrire $t$ est son \textit{rang}. Montrer que pour
$t \in V \otimes W$, on retrouve ainsi le rang usuel d'une matrice.
\item[(c)] Montrer que la multiplication des matrices $2 \times 2$
\og{} est \fg{} un tenseur précis $T \in U \otimes V \otimes W$ avec $\dim U = \dim V = \dim W = 4$, et
que l'algorithme de Strassen (LIN-17) est exactement une décomposition de ce tenseur en 7 tenseurs
purs.
\item[(d)] Contrairement au rang matriciel, le rang tensoriel à trois facteurs peut sauter \textit{vers le bas}
à la limite : exhiber une suite de tenseurs de rang 2 convergeant vers un tenseur de rang 3. (Ce
phénomène --- le rang de bord --- est la raison pour laquelle la complexité de la multiplication
matricielle est si difficile à cerner.)
\end{enumerate}

\end{problem}

\noindent Les produits tensoriels sont entrés en mathématiques par la théorie des invariants
et la géométrie différentielle du dix-neuvième siècle, ont été faits algèbre par Whitney (1938), et
rendus incontournables par la mécanique quantique, où l'espace des états de deux systèmes est le
produit tensoriel de leurs espaces d'états --- l'intrication est exactement l'existence de tenseurs non
purs, partie (ii). La partie (iii), l'observation que les algorithmes rapides sont des décompositions
tensorielles de rang faible, est due à Strassen et a transformé une question algorithmique en
géométrie algébrique ; le rang de bord de la partie (iv) (Bini et al., 1979) en est la subtilité
centrale. Ce problème est la porte de l'algèbre linéaire vers l'algèbre \textit{multilinéaire}.

\begin{refs}
\item Contexte : Produit tensoriel --- Wikipédia (\url{https://fr.wikipedia.org/wiki/Produit_tensoriel})
\item Contexte : Décomposition de rang tensoriel --- Wikipédia (en anglais) (\url{https://en.wikipedia.org/wiki/Tensor_rank_decomposition})
\item Lecture : J. Landsberg, \textit{Tensors: Geometry and Applications} (AMS GSM 128), ch. 1--2
\end{refs}

\bigskip

\begin{problem}[{La nilpotence détectée par les traces --- Master, short}]
\textbf{LIN-28.} \textbf{Problème.} Soit $N$ une matrice $n \times n$ sur $\mathbb{C}$.
Démontrer que les assertions suivantes sont équivalentes : $N^k = 0$ pour un certain $k$ ; toute
valeur propre de $N$ est $0$ ; et $\operatorname{tr}(N^k) = 0$ pour $k = 1, 2, \dots, n$.
\end{problem}

\noindent La troisième condition est remarquable parce qu'il s'agit d'une liste finie de
conditions linéaires sur les coefficients de $N^k$, détectant une propriété purement qualitative. Le
pont est constitué par les identités de Newton, qui expriment les coefficients du polynôme
caractéristique au moyen des sommes de puissances $\operatorname{tr}(N^k)$ --- et c'est exactement là
qu'intervient l'hypothèse de caractéristique nulle : sur $\mathbb{F}_p$, la matrice identité de
taille $p$ a toutes ses traces de puissances nulles sans être nilpotente. L'énoncé est l'ombre
matricielle du théorème d'Engel en théorie de Lie.

\begin{refs}
\item Contexte : Matrice nilpotente --- Wikipédia (\url{https://fr.wikipedia.org/wiki/Matrice_nilpotente})
\item Contexte : Identités de Newton --- Wikipédia (\url{https://fr.wikipedia.org/wiki/Identit%C3%A9s_de_Newton})
\item Lecture : R. Horn \& C. Johnson, \textit{Matrix Analysis}, 2e éd., ch. 1--2
\end{refs}

\bigskip

\begin{problem}[{Une seule base pour deux matrices --- Master, short}]
\textbf{LIN-29.} \textbf{Problème.} Soient $A$ et $B$ deux matrices complexes $n \times n$
diagonalisables. Démontrer qu'il existe une seule et même matrice inversible $S$ rendant
$S^{-1} A S$ et $S^{-1} B S$ toutes deux diagonales si et seulement si $AB = BA$, et montrer par un exemple que la
conclusion tombe en défaut si l'on suppose seulement que les matrices commutent sans être
diagonalisables.
\end{problem}

\noindent C'est le contenu algébrique d'un slogan physique : deux observables peuvent être
mesurées simultanément précisément lorsque leurs opérateurs commutent, car c'est seulement alors
qu'elles partagent une base propre. La démonstration --- restreindre l'un des opérateurs aux espaces
propres de l'autre --- est le prototype du théorème spectral pour les familles commutatives d'opérateurs
normaux, et de la diagonalisation simultanée d'une algèbre commutative entière, dont les matrices
circulantes de LIN-30 sont l'exemple le plus net.

\begin{refs}
\item Contexte : Paire de matrices commutantes --- Wikipédia (\url{https://fr.wikipedia.org/wiki/Paire_de_matrices_commutantes})
\item Contexte : Matrice diagonalisable --- Wikipédia (\url{https://fr.wikipedia.org/wiki/Matrice_diagonalisable})
\item Lecture : R. Horn \& C. Johnson, \textit{Matrix Analysis}, 2e éd., ch. 1.3
\end{refs}

\bigskip

\begin{problem}[{Les matrices circulantes, et pourquoi la matrice de Fourier apparaît --- Master}]
\textbf{LIN-30.} \textbf{Problème.} Une \textit{circulante} est une matrice $C$ dont les lignes sont les décalages
cycliques successifs de la première : son coefficient $(j,k)$ est $c_{k-j}$, les indices étant lus
modulo $n$. Soit $S$ la circulante de première ligne $(0, 1, 0, \dots, 0)$ --- le décalage
cyclique --- et posons $\omega = e^{2\pi i / n}$.
\begin{enumerate}
\item[(a)] Démontrer que $S^n = I$ et que toute circulante est égale à $p(S)$ pour le polynôme
$p(z) = c_0 + c_1 z + \cdots + c_{n-1} z^{n-1}$. Démontrer réciproquement qu'une matrice est
circulante si et seulement si elle commute avec $S$.
\item[(b)] Démontrer que les vecteurs $v_k = (1, \omega^{k}, \omega^{2k}, \dots, \omega^{(n-1)k})^{\mathsf{T}}$,
pour $k = 0, 1, \dots, n-1$, sont vecteurs propres de \textit{toute} circulante $n \times n$. En
déduire que la matrice unitaire $F$ de coefficients $F_{jk} = n^{-1/2} \omega^{jk}$ --- la
transformation de Fourier discrète --- diagonalise simultanément toutes les circulantes, avec
\[ F^{*} C F = \operatorname{diag}\big(p(1), p(\omega), \dots, p(\omega^{n-1})\big). \]
\item[(c)] En déduire que les circulantes forment une algèbre commutative isomorphe à
$\mathbb{C}[z] / (z^n - 1)$, que $\det C = \prod_{k=0}^{n-1} p(\omega^{k})$, et que $Cx$ est la
convolution cyclique de $c$ avec $x$ --- de sorte qu'elle se calcule en $O(n \log n)$ opérations
par transformation de Fourier rapide.
\item[(d)] Démontrer que toute circulante est une matrice normale, et donner une formule pour $C^{-1}$
lorsqu'elle existe. Expliquer ensuite pourquoi la partie (ii) est un cas particulier d'un principe
général : une application linéaire commutant avec l'action d'un groupe abélien fini est diagonalisée
par les caractères de ce groupe.
\end{enumerate}

\end{problem}

\noindent La formule des valeurs propres de (ii) fut trouvée au dix-neuvième siècle --- Catalan
et Muir ont tous deux calculé les déterminants circulants --- mais sa signification relève de la théorie
des représentations : une circulante est exactement un opérateur commutant avec le décalage cyclique,
et la base de Fourier discrète est la base des caractères de $\mathbb{Z}/n\mathbb{Z}$. Cette seule
observation explique pourquoi la transformation de Fourier change la convolution en multiplication, et
pourquoi la FFT de Cooley--Tukey de 1965 --- anticipée par Gauss dans une note non publiée de 1805 --- est
l'algorithme le plus utilisé en ingénierie. C'est aussi le modèle fini de la façon dont la symétrie
diagonalise les opérateurs différentiels : les ondes planes sont fonctions propres du laplacien pour
la même raison.

\begin{refs}
\item Contexte : Matrice circulante --- Wikipédia (\url{https://fr.wikipedia.org/wiki/Matrice_circulante})
\item Contexte : Transformation de Fourier discrète --- Wikipédia (\url{https://fr.wikipedia.org/wiki/Transformation_de_Fourier_discr%C3%A8te})
\item Lecture : P. J. Davis, \textit{Circulant Matrices}
\item Cours : MIT OCW 18.065 Matrix Methods in Data Analysis (en anglais) (\url{https://ocw.mit.edu/courses/18-065-matrix-methods-in-data-analysis-signal-processing-and-machine-learning-spring-2018/})
\end{refs}

\bigskip

\begin{problem}[{Le triangle de Schur, et ce que \og{} normal \fg{} veut vraiment dire --- Master}]
\textbf{LIN-39.} \textbf{Problème.} Soit $A$ une matrice complexe $n \times n$ et $A^{*}$ sa transconjuguée. On
dit que $A$ est \textit{normale} si $A A^{*} = A^{*} A$, et l'on note $\|A\|_F$ la norme de
Frobenius $\big(\sum_{i,j} |a_{ij}|^2\big)^{1/2}$.
\begin{enumerate}
\item[(a)] Démontrer le théorème de Schur : $A = Q T Q^{*}$ avec $Q$ unitaire et $T$ triangulaire
supérieure, et les valeurs propres peuvent être amenées à apparaître le long de la diagonale de $T$
dans n'importe quel ordre prescrit.
\item[(b)] En déduire que $A$ est normale si et seulement si elle est unitairement diagonalisable. En
dériver les théorèmes spectraux pour les matrices hermitiennes, antihermitiennes et unitaires, et dire
dans chaque cas exactement où dans $\mathbb{C}$ les valeurs propres sont contraintes de se
trouver.
\item[(c)] Démontrer l'inégalité de Schur : si $\lambda_1, \dots, \lambda_n$ sont les valeurs propres de
$A$ comptées avec multiplicité, alors $\sum_i |\lambda_i|^2 \le \|A\|_F^2$, avec égalité si et
seulement si $A$ est normale.
\item[(d)] Démontrer deux caractérisations supplémentaires : $A$ est normale si et seulement si
$\|Ax\| = \|A^{*}x\|$ pour tout $x$, et si et seulement si $A^{*} = p(A)$ pour un certain
polynôme $p$. Démontrer enfin la forme de Schur réelle : toute matrice carrée réelle est
orthogonalement semblable à une matrice triangulaire supérieure par blocs dont les blocs diagonaux
sont $1 \times 1$ et $2 \times 2$.
\end{enumerate}

\end{problem}

\noindent Issai Schur publia la triangularisation en 1909. Élève de Frobenius et principal
algébriste de Berlin, il fut chassé de sa chaire par les lois raciales nazies, émigra en Palestine en
1939 et mourut à Tel-Aviv deux ans plus tard. Son théorème est ce que la réduction de Jordan de
LIN-14 ne peut pas être : le changement de base est unitaire, donc il ne peut amplifier les erreurs,
alors que la similitude produisant une réduite de Jordan peut être arbitrairement mal conditionnée et
que la structure de Jordan elle-même est détruite par une perturbation arbitrairement petite. C'est
pourquoi l'algorithme QR calcule une forme de Schur et non une forme de Jordan, et pourquoi la
normalité est \textit{la} ligne de partage de l'algèbre linéaire numérique : pour les matrices normales,
les valeurs propres sont parfaitement conditionnées, et tout ce qui est difficile --- pseudospectres,
croissance transitoire, conjecture de Crouzeix (LIN-21) --- vit strictement de l'autre côté.

\begin{refs}
\item Contexte : Décomposition de Schur --- Wikipédia (\url{https://fr.wikipedia.org/wiki/D%C3%A9composition_de_Schur})
\item Contexte : Matrice normale --- Wikipédia (\url{https://fr.wikipedia.org/wiki/Matrice_normale})
\item Lecture : R. Horn \& C. Johnson, \textit{Matrix Analysis}, 2e éd., ch. 2
\item Lecture : L. Trefethen \& D. Bau, \textit{Numerical Linear Algebra}, partie V
\end{refs}

\bigskip

\begin{problem}[{Produits de Kronecker et équation de Sylvester --- Master}]
\textbf{LIN-40.} \textbf{Problème.} Pour $A$ de taille $m \times m$ et $B$ de taille $n \times n$, le
\textit{produit de Kronecker} $A \otimes B$ est la matrice $mn \times mn$ formée des blocs
$a_{ij}B$, et $\operatorname{vec}(X)$ désigne le long vecteur obtenu en empilant les colonnes de
$X$.
\begin{enumerate}
\item[(a)] Démontrer la règle du produit mixte $(A \otimes B)(C \otimes D) = (AC) \otimes (BD)$ chaque fois
que les produits sont définis, et en déduire que $A \otimes B$ est inversible exactement lorsque
$A$ et $B$ le sont toutes deux.
\item[(b)] Démontrer que les valeurs propres de $A \otimes B$ sont exactement les produits
$\lambda_i \mu_j$ des valeurs propres de $A$ et de $B$, comptées avec multiplicité, et en
déduire $\operatorname{tr}(A \otimes B) = \operatorname{tr}(A)\operatorname{tr}(B)$ et
$\det(A \otimes B) = (\det A)^{n} (\det B)^{m}$.
\item[(c)] Démontrer l'identité
\[ \operatorname{vec}(AXB) = \big(B^{\mathsf{T}} \otimes A\big)\operatorname{vec}(X), \]
et l'utiliser, avec (b), pour démontrer le théorème de Sylvester : l'équation $AX - XB = C$ admet une
unique solution $X$ pour tout $C$ si et seulement si $A$ et $B$ n'ont aucune valeur propre en
commun.
\item[(d)] En déduire le cas de Lyapunov : si toute valeur propre de $A$ est de partie réelle strictement
négative, alors $A^{\mathsf{T}}P + PA = -Q$ admet exactement une solution $P$ pour chaque $Q$, et
$P$ est symétrique définie positive dès que $Q$ l'est. Expliquer ce que cela dit du comportement à
long terme de $x' = Ax$ (LIN-12), et pourquoi cela transforme la \og{} stabilité \fg{} en un test de
faisabilité semi-défini.
\end{enumerate}

\end{problem}

\noindent Le produit porte le nom de Kronecker sur de minces indices : Johann Georg Zehfuss
l'écrivit en 1858, comme l'ont documenté Henderson, Pukelsheim et Searle dans une étude historique de
1983, et il circula pendant des décennies sous le nom de \og{} matrice de Zehfuss \fg{}. Sylvester posa son
équation en 1884, et le critère spectral de (c) a été redécouvert à plusieurs reprises depuis, y
compris sous forme d'opérateurs par Rosenblum en 1956. Conceptuellement, ce problème est LIN-18 en
coordonnées : $A \otimes B$ est la matrice du produit tensoriel de deux applications linéaires, et
la partie (b) énonce que des systèmes indépendants multiplient leurs spectres --- ce qui est exactement
pourquoi l'espace des états de deux systèmes quantiques est un produit tensoriel. Pratiquement,
l'équation de Lyapunov de (d) est la raison pour laquelle les automaticiens peuvent certifier la
stabilité en résolvant un seul système linéaire plutôt qu'en intégrant une équation
différentielle.

\begin{refs}
\item Contexte : Produit de Kronecker --- Wikipédia (\url{https://fr.wikipedia.org/wiki/Produit_de_Kronecker})
\item Contexte : Équation de Sylvester --- Wikipédia (\url{https://fr.wikipedia.org/wiki/%C3%89quation_de_Sylvester})
\item Article : H. V. Henderson, F. Pukelsheim \& S. R. Searle, \og{} On the history of the Kronecker product \fg{} (1983), Linear and Multilinear Algebra 14, 113--120
\item Lecture : R. Horn \& C. Johnson, \textit{Topics in Matrix Analysis}, ch. 4
\end{refs}

\bigskip

\begin{problem}[{Le théorème de Birkhoff sur les matrices bistochastiques --- Master, short}]
\textbf{LIN-41.} \textbf{Problème.} On dit qu'une matrice $n \times n$ est \textit{bistochastique}
si ses coefficients sont positifs ou nuls et si la somme de chaque ligne et de chaque colonne vaut
$1$. Démontrer le théorème de Birkhoff--von Neumann : ces matrices forment un convexe compact dont
les points extrémaux sont exactement les $n!$ matrices de permutation, de sorte que toute matrice
bistochastique est une combinaison convexe de matrices de permutation.
\end{problem}

\noindent Birkhoff publia ce résultat en 1946, en trois brèves observations dans une revue de
l'université de Tucumán ; von Neumann démontra un énoncé équivalent en 1953 en travaillant sur le
problème d'affectation. Le théorème est en réalité le théorème des mariages de Hall en costume
matriciel --- le nœud de toute démonstration consiste à extraire une permutation du motif des
coefficients strictement positifs --- et il dit que les \og{} permutations randomisées \fg{} sont le seul
comportement bistochastique qui existe. Le polytope qu'il décrit est de dimension $(n-1)^2$ et
possède $n!$ sommets, et son volume n'est connu que pour de petits $n$ ; la question naturelle
suivante, celle de savoir à quel point le permanent d'une matrice bistochastique peut être petit,
était la conjecture de van der Waerden de 1926, réglée indépendamment par Egorychev et Falikman en
1980--81.

\begin{refs}
\item Contexte : Matrice bistochastique --- Wikipédia (\url{https://fr.wikipedia.org/wiki/Matrice_bistochastique})
\item Contexte : Polytope de Birkhoff --- Wikipédia (en anglais) (\url{https://en.wikipedia.org/wiki/Birkhoff_polytope})
\item Lecture : R. A. Brualdi, \textit{Combinatorial Matrix Classes}, ch. 9
\end{refs}

\bigskip


\section*{Recherche}

\begin{problem}[{L'exposant de la multiplication matricielle --- Recherche}]
\textbf{LIN-19.} \textbf{Problème.} L'exposant $\omega$ est la borne inférieure de tous les $c$
tels que deux matrices $n \times n$ puissent être multipliées en $O(n^c)$ opérations
arithmétiques. A-t-on $\omega = 2$ ? \textbf{Statut : ouvert.} Ce que l'on sait :
$2 \le \omega \le \log_2 7 \approx 2{,}807$ (Strassen 1969, LIN-17) ; $\omega < 2{,}376$
(Coppersmith--Winograd, 1990) ; après un long plateau, $\omega < 2{,}3729$ (Le Gall 2014 ;
Alman--Vassilevska Williams 2021), $\omega < 2{,}37188$ (Duan--Wu--Zhou, 2023),
$\omega < 2{,}371552$ (Vassilevska Williams--Xu--Xu--Zhou, SODA 2024), $\omega < 2{,}371339$
(Alman--Duan--Vassilevska Williams--Xu--Xu--Zhou, SODA 2025). En août 2026, le même cadre de la méthode
laser, avec une chaîne d'optimisation assistée par le système AlphaEvolve, a atteint
$\omega < 2{,}371177$. Des barrières connues démontrent que le tenseur de Coppersmith--Winograd ne
peut mener ces méthodes en dessous d'environ 2,3 : combler l'écart jusqu'à 2 --- ou démontrer
$\omega > 2$ --- demande une idée véritablement neuve. Pour entrer dans le sujet : démontrez
vous-même $\omega \le 2{,}81$ (LIN-17), puis lisez l'approche par la théorie des groupes de Cohn et
Umans, qui donnerait $\omega = 2$ si certaines familles de groupes existent.
\end{problem}

\noindent Savoir si le calcul le plus fondamental de l'algèbre linéaire peut se faire en
temps essentiellement quadratique compte parmi les grands problèmes ouverts de l'informatique et de
l'algèbre à la fois, relié par LIN-18 à la géométrie d'un seul tenseur. Les progrès depuis 2010 se
font par incréments de $10^{-3}$ puis de $10^{-4}$ --- chaque incrément étant un article
substantiel --- et l'étape de 2026 s'est fait remarquer comme une collaboration entre humain et
machine. De nombreux théoriciens de la complexité croient que $\omega = 2$ ; aucun ne sait le
démontrer, et les algorithmes rapides praticables restent très en retrait de ces bornes théoriques,
qui dissimulent des constantes astronomiques.

\begin{refs}
\item Contexte : Complexité de la multiplication de matrices --- Wikipédia (\url{https://fr.wikipedia.org/wiki/Complexit%C3%A9_de_la_multiplication_de_matrices})
\item Article : V. Vassilevska Williams, Y. Xu, Z. Xu, R. Zhou, \og{} New Bounds for Matrix Multiplication: from Alpha to Omega \fg{} (2024), arXiv:2307.07970 (\url{https://arxiv.org/abs/2307.07970})
\item Article : J. Alman, R. Duan, V. Vassilevska Williams, Y. Xu, Z. Xu, R. Zhou, \og{} More Asymmetry Yields Faster Matrix Multiplication \fg{} (2025), arXiv:2404.16349 (\url{https://arxiv.org/abs/2404.16349})
\item Article : E. Dupont et al., \og{} Improving the matrix multiplication exponent with modern optimization and AlphaEvolve \fg{} (2026), arXiv:2608.16884 (\url{https://arxiv.org/abs/2608.16884})
\end{refs}

\bigskip

\begin{problem}[{La conjecture de Hadamard --- Recherche}]
\textbf{LIN-20.} \textbf{Problème.} Une \textit{matrice de Hadamard} d'ordre $n$ est une
matrice $n \times n$ à coefficients $\pm 1$ dont les lignes sont deux à deux orthogonales. La
conjecture de Hadamard : une telle matrice existe pour tout $n$ divisible par 4.
\textbf{Statut : ouvert.} Pour entrer dans le sujet :
\begin{enumerate}
\item[(a)] démontrer qu'une matrice de Hadamard d'ordre $n > 2$ force $n \equiv 0 \pmod 4$ ;
\item[(b)] démontrer l'inégalité de Hadamard --- une matrice $n \times n$ à coefficients dans $[-1, 1]$
vérifie $|\det| \le n^{n/2}$ --- et montrer que l'égalité a lieu exactement pour les matrices de
Hadamard ;
\item[(c)] construire des matrices de Hadamard de tous les ordres $2^k$ (Sylvester, 1867) et d'ordre 12
(construction de Paley, 1933). Ce que l'on sait : les constructions couvrent la plupart des ordres ;
le plus petit cas indécis se situait à 428 jusqu'à ce que Kharaghani et Tayfeh-Rezaie en trouvent une
en 2004, puis à 668 pendant deux décennies --- jusqu'en août 2026, où une matrice d'ordre 668 a été
annoncée par Levent Alpöge et ses collaborateurs (un certificat vérifiable par machine, trouvé avec
assistance informatique). Les plus petits ordres encore indécis sont 716 et 892 ; la conjecture
générale reste grande ouverte.
\end{enumerate}

\end{problem}

\noindent Les matrices de Hadamard sont entrées par l'inégalité sur le déterminant (Hadamard,
1893), mais Sylvester avait joué avec la famille des $2^k$ des décennies plus tôt. Elles relèvent à
la fois de la combinatoire extrémale, de la théorie des codes (le code de Reed--Muller embarqué sur
Mariner 9) et des plans d'expérience (Plackett--Burman). La forme de la conjecture est typique de la
combinatoire difficile : une densité écrasante d'exemples, aucune construction générale --- chaque ordre
récalcitrant n'est tombé que devant une idée algébrique sur mesure, plus, dans le dernier cas, un
calcul sérieux.

\begin{refs}
\item Contexte : Matrice de Hadamard --- Wikipédia (\url{https://fr.wikipedia.org/wiki/Matrice_de_Hadamard})
\item Contexte : Construction de Paley --- Wikipédia (en anglais) (\url{https://en.wikipedia.org/wiki/Paley_construction})
\item Article : H. Kharaghani \& B. Tayfeh-Rezaie, \og{} A Hadamard matrix of order 428 \fg{} (2005), Journal of Combinatorial Designs 13, 435--440
\end{refs}

\bigskip

\begin{problem}[{La conjecture de Crouzeix --- Recherche}]
\textbf{LIN-21.} \textbf{Problème.} L'\textit{image numérique} d'une matrice carrée complexe
$A$ est l'ensemble $W(A) = \{ x^* A x : \|x\| = 1 \}$ du plan complexe --- un convexe compact
contenant les valeurs propres. Conjecture de Crouzeix (2004) : pour toute matrice $A$ et tout
polynôme $p$,
\[ \| p(A) \| \le 2 \max_{z \in W(A)} |p(z)|, \]
où $\|\cdot\|$ est la norme d'opérateur. \textbf{Statut : ouvert.} Ce que l'on sait : l'inégalité vaut
avec la constante 11,08 (Crouzeix, 2007) et avec la constante $1 + \sqrt{2}$ (Crouzeix--Palencia,
2017) ; la constante 2 serait optimale, et la conjecture est vérifiée pour les matrices
$2 \times 2$ et diverses classes particulières. Pour entrer dans le sujet :
\begin{enumerate}
\item[(a)] démontrer que $W(A)$ est convexe (le
théorème de Toeplitz--Hausdorff) ;
\item[(b)] démontrer la conjecture --- avec la constante 1 --- pour les matrices normales,
en utilisant le théorème spectral ;
\item[(c)] montrer par un exemple que pour $A$ non normale, $\|p(A)\|$ peut
dépasser $\max_{W(A)} |p|$, de sorte qu'une constante strictement supérieure à 1 est réellement
nécessaire.
\end{enumerate}

\end{problem}

\noindent La conjecture demande jusqu'où une matrice non normale peut amplifier un polynôme
au-delà de ce que prédit son \og{} ombre \fg{} dans le plan --- la forme quantitative la plus serrée de l'écart
entre comportement normal et non normal qui domine l'analyse numérique (le même écart qui se cache
derrière les pseudospectres et la croissance transitoire). Michel Crouzeix la posa après avoir démontré
la première constante finie ; l'élégante démonstration de $1 + \sqrt 2$ par Crouzeix et Palencia en
2017 a ranimé le domaine, et le dernier pas de $1+\sqrt{2}$ à 2 a résisté à une décennie d'efforts,
tant des théoriciens des opérateurs que des analystes numériciens.

\begin{refs}
\item Contexte : Conjecture de Crouzeix --- Wikipédia (en anglais) (\url{https://en.wikipedia.org/wiki/Crouzeix%27s_conjecture})
\item Contexte : Image numérique --- Wikipédia (en anglais) (\url{https://en.wikipedia.org/wiki/Numerical_range})
\item Article : M. Crouzeix \& C. Palencia, \og{} The numerical range is a $(1+\sqrt{2})$-spectral set \fg{} (2017), SIAM J. Matrix Anal. Appl. 38, 649--655
\item Voir aussi : FA-29 (\url{pure-functional-analysis.html#FA-29}) donne la formulation en termes d'opérateurs sur un espace de Hilbert de la même conjecture.
\end{refs}

\bigskip

\begin{problem}[{La conjecture de Hadamard circulante --- Recherche, short}]
\textbf{LIN-31.} \textbf{Problème.} Une matrice de Hadamard circulante est une circulante (LIN-30) à
coefficients $\pm 1$ dont les lignes sont deux à deux orthogonales (LIN-20) ; les seuls exemples
connus sont d'ordre $1$ et $4$. Démontrer qu'il n'en existe pas d'autres, ou en exhiber une d'ordre
strictement supérieur à $4$ --- \textbf{statut : ouvert.}
\end{problem}

\noindent Ryser conjectura cela vers 1963, et c'est l'un de ces rares problèmes ouverts où le
premier pas est un véritable plaisir : diagonaliser par la matrice de Fourier (LIN-30) transforme
l'orthogonalité en la condition $|p(\omega^k)| = \sqrt{n}$ pour toute racine $n$-ième de l'unité,
ce qui force $n$ à valoir quatre fois un carré parfait et livre le problème à la théorie algébrique
des nombres. Les travaux de Turyn sur les nombres premiers autoconjugués et la méthode de la
\og{} descente de corps \fg{} de Leung et Schmidt ont depuis éliminé tous les ordres $4 < n \le 10^{11}$
sauf une poignée de candidats --- un domaine vérifié considérable, sans démonstration en vue. La
conjecture équivaut à l'inexistence de certaines suites binaires parfaites, ce qui explique pourquoi
les ingénieurs du radar et des télécommunications s'y intéressent.

\begin{refs}
\item Contexte : Matrice circulante --- Wikipédia (\url{https://fr.wikipedia.org/wiki/Matrice_circulante})
\item Contexte : Matrice de Hadamard --- Wikipédia (\url{https://fr.wikipedia.org/wiki/Matrice_de_Hadamard})
\item Article : K. H. Leung \& B. Schmidt, \og{} New restrictions on possible orders of circulant Hadamard matrices \fg{} (2012), Designs, Codes and Cryptography
\end{refs}

\bigskip

\begin{problem}[{Permanent contre déterminant --- Recherche, short}]
\textbf{LIN-32.} \textbf{Problème.} Le permanent d'une matrice $n \times n$ est
$\operatorname{per}(X) = \sum_{\sigma \in S_n} \prod_{i=1}^{n} x_{i \sigma(i)}$ --- le déterminant
privé de tous ses signes moins --- et sa \textit{complexité déterminantale} $\mathrm{dc}(n)$ est le plus
petit $m$ pour lequel $\operatorname{per}(X) = \det(M)$ pour une certaine matrice $m \times m$
$M$ dont les coefficients sont des fonctions affines des coefficients de $X$. Déterminer la
croissance de $\mathrm{dc}(n)$ : Valiant a conjecturé qu'elle est superpolynomiale, et
\textbf{statut : ouvert} --- les meilleures bornes connues sont $\mathrm{dc}(n) \ge n^2/2$
(Mignon--Ressayre, 2004) contre une borne supérieure exponentielle d'environ $2^n$.
\end{problem}

\noindent Valiant montra en 1979 que le permanent est complet pour la classe de comptage
$\#\mathrm{P}$ alors que le déterminant se calcule en temps polynomial, de sorte que ces deux
polynômes presque identiques siègent de part et d'autre d'une frontière de complexité ; les séparer
est l'analogue algébrique de P contre NP. Le problème s'énonce entièrement en termes d'algèbre
linéaire, et il est à l'origine de la théorie géométrique de la complexité, le programme de Mulmuley
et Sohoni pour l'attaquer par la théorie des représentations et la géométrie algébrique. L'écart entre
$n^2/2$ et $2^n$ n'a presque pas bougé en deux décennies.

\begin{refs}
\item Contexte : Permanent (mathématiques) --- Wikipédia (\url{https://fr.wikipedia.org/wiki/Permanent_%28math%C3%A9matiques%29})
\item Contexte : Calcul du permanent --- Wikipédia (en anglais) (\url{https://en.wikipedia.org/wiki/Computing_the_permanent})
\item Article : T. Mignon \& N. Ressayre, \og{} A quadratic bound for the determinant and permanent problem \fg{} (2004), International Mathematics Research Notices 2004, no 79, 4241--4253
\end{refs}

\bigskip

\begin{problem}[{Combien de droites équiangulaires tiennent dans l'espace ? --- Recherche}]
\textbf{LIN-33.} \textbf{Problème.} Une famille de droites passant par l'origine de $\mathbb{R}^d$ est
\textit{équiangulaire} si deux quelconques d'entre elles font le même angle. Notons $N(d)$ le plus
grand nombre de telles droites, et $N_\alpha(d)$ le plus grand nombre lorsque l'angle commun
$\theta$ est fixé avec $\alpha = \cos\theta$.
\begin{enumerate}
\item[(a)] Démontrer la borne de Gerzon $N(d) \le \tbinom{d+1}{2}$, en montrant que les projections
orthogonales sur les droites sont linéairement indépendantes à l'intérieur de l'espace de dimension
$\tbinom{d+1}{2}$ des matrices symétriques $d \times d$.
\item[(b)] Démontrer la borne relative : si les droites sont engendrées par des vecteurs unitaires $u_i$
vérifiant $|\langle u_i, u_j \rangle| = \alpha$ pour $i \ne j$, et si $d < 1/\alpha^2$, alors
\[ N_\alpha(d) \;\le\; \frac{d\,(1 - \alpha^2)}{1 - d\,\alpha^2}. \]
\item[(c)] Construire 6 droites équiangulaires dans $\mathbb{R}^3$ (les six diagonales d'un icosaèdre
régulier) et 28 dans $\mathbb{R}^7$, et démontrer que 6 est optimal dans $\mathbb{R}^3$.
\item[(d)] Étudier le théorème de 2021 de Jiang, Tidor, Yao, Zhang et Zhao, qui détermine $N_\alpha(d)$
pour tout $\alpha$ fixé et tout $d$ suffisamment grand, et rédiger un exposé de son ingrédient
spectral clé : un graphe connexe de degré borné a une multiplicité sous-linéaire de sa deuxième valeur
propre. \textbf{Statut :} le problème à angle fixé est réglé pour $d$ grand ; la valeur exacte de
$N(d)$ n'est connue que pour une courte liste de dimensions et reste ouverte en général, tout comme
le comportement lorsque l'on autorise l'angle à décroître avec $d$.
\end{enumerate}

\end{problem}

\noindent Haantjes demanda en 1948 combien de droites équiangulaires tiennent dans l'espace de
dimension trois ; van Lint et Seidel entamèrent la théorie systématique en 1966, et le mémoire de
Lemmens et Seidel de 1973 en fit une famille de conjectures précises qui tinrent près de cinquante
ans. Le sujet est un carrefour : les systèmes de droites équiangulaires sont les mêmes données que
certains bigraphes et graphes fortement réguliers, ce sont des codes sphériques optimaux, et leur
analogue complexe --- $d^2$ droites équiangulaires dans $\mathbb{C}^d$, les SIC-POVM de
l'information quantique --- est lui-même un fameux problème ouvert. La résolution de 2021 du cas à angle
fixé est venue d'une direction inattendue, la théorie spectrale des graphes, et le régime exponentiel
est aujourd'hui activement attaqué.

\begin{refs}
\item Contexte : Droites équiangulaires --- Wikipédia (en anglais) (\url{https://en.wikipedia.org/wiki/Equiangular_lines})
\item Article : Z. Jiang, J. Tidor, Y. Yao, S. Zhang, Y. Zhao, \og{} Equiangular lines with a fixed angle \fg{} (2021), Annals of Mathematics 194, 729--743, arXiv:1907.12466 (\url{https://arxiv.org/abs/1907.12466})
\item Article : I. Balla \& M. Bucić, \og{} Equiangular lines via improved eigenvalue multiplicity \fg{} (2024), arXiv:2409.16219 (\url{https://arxiv.org/abs/2409.16219})
\end{refs}

\bigskip

\begin{problem}[{À quelle vitesse un produit de matrices peut-il croître ? --- Recherche}]
\textbf{LIN-42.} \textbf{Problème.} Pour un ensemble fini $\mathcal{M} = \{A_1, \dots, A_m\}$ de matrices réelles
$n \times n$, le \textit{rayon spectral joint} est
\[ \rho(\mathcal{M}) \;=\; \lim_{k \to \infty} \; \max \big\{ \|A_{i_1} A_{i_2} \cdots A_{i_k}\|^{1/k} \;:\; A_{i_j} \in \mathcal{M} \big\}, \]
le taux de croissance asymptotique le plus rapide réalisable par des produits pris dans
$\mathcal{M}$.
\begin{enumerate}
\item[(a)] Démontrer que la limite existe et ne dépend pas de la norme sous-multiplicative employée (lemme
sous-additif de Fekete). Pour $m = 1$, démontrer qu'il s'agit du rayon spectral ordinaire --- la
formule de Gelfand.
\item[(b)] Démontrer que tout produit de matrices de $\mathcal{M}$ tend vers $0$ lorsque sa longueur
croît si et seulement si $\rho(\mathcal{M}) < 1$, et démontrer l'inégalité
$\rho(\mathcal{M}) \ge \rho(A_{i_1} \cdots A_{i_k})^{1/k}$ pour tout produit fini. (Que la borne
supérieure du membre de droite soit égale à $\rho$ est le théorème de Berger--Wang, 1992.)
\item[(c)] Démontrer que $\rho(\mathcal{M}) = \max_i \rho(A_i)$ lorsque les matrices sont simultanément
trigonalisables, et exhiber un couple de matrices $2 \times 2$ pour lequel
$\rho(\mathcal{M}) > \max_i \rho(A_i)$ --- chaque matrice prise séparément peut donc être stable
alors que l'un de leurs produits ne l'est pas.
\item[(d)] Lire l'état du sujet et en rédiger un exposé. \textbf{Statut :} la conjecture de finitude de Lagarias
et Wang (1995) --- selon laquelle $\rho$ est toujours atteint par un produit périodique --- est
\textbf{fausse} (Bousch et Mairesse, 2002 ; un contre-exemple explicite a été donné par Hare, Morris,
Sidorov et Theys en 2011). Calculer ou approcher $\rho$ est NP-difficile, et décider si
$\rho \le 1$ est \textbf{indécidable} (Blondel et Tsitsiklis). Savoir si la question stricte \og{} a-t-on
$\rho < 1$ ? \fg{} est décidable est, autant qu'on sache, \textbf{ouvert} --- de même que la description
des ensembles de matrices qui possèdent effectivement la propriété de finitude. En attendant, des
algorithmes pratiques approchent bien $\rho$ pour les ensembles rencontrés dans les
applications.
\end{enumerate}

\end{problem}

\noindent Rota et Strang introduisirent cette quantité au MIT en 1960 et elle resta discrète
jusqu'à ce que Daubechies et Lagarias découvrent, à la fin des années 1980, que la régularité d'une
ondelette est gouvernée par le rayon spectral joint de deux matrices construites à partir de ses
coefficients de raffinement. Elle décide aujourd'hui de la stabilité des systèmes commutés en
automatique, de la convergence des schémas de subdivision en conception assistée par ordinateur, et de
la croissance de suites combinatoires. La situation mathématique est d'une netteté inhabituelle : la
quantité est définie par une limite d'une seule ligne, elle est calculable en pratique avec la
précision que l'on veut, et pourtant les questions naturelles à son sujet sont prouvablement
intraitables ou indécidables --- un point de rencontre concret entre l'algèbre linéaire et la
calculabilité, qui récompense une lecture lente.

\begin{refs}
\item Contexte : Rayon spectral joint --- Wikipédia (en anglais) (\url{https://en.wikipedia.org/wiki/Joint_spectral_radius})
\item Lecture : R. Jungers, \textit{The Joint Spectral Radius: Theory and Applications} (Springer, 2009)
\item Article : T. Bousch \& J. Mairesse, \og{} Asymptotic height optimization for topical IFS, Tetris heaps, and the finiteness conjecture \fg{} (2002), Journal of the American Mathematical Society
\item Article : K. G. Hare, I. D. Morris, N. Sidorov \& J. Theys, \og{} An explicit counterexample to the Lagarias--Wang finiteness conjecture \fg{} (2011), Advances in Mathematics
\end{refs}

\bigskip

\begin{problem}[{La conjecture BMV, aujourd'hui théorème de Stahl --- Recherche, short}]
\textbf{LIN-43.} \textbf{Problème.} Soient $A$ et $B$ des matrices hermitiennes
$n \times n$ semi-définies positives et soit $m \ge 0$ un entier ; démontrer que tout coefficient
du polynôme $t \mapsto \operatorname{tr}\big((A + tB)^{m}\big)$ est positif ou nul.
\textbf{Statut : démontré} --- c'est la forme de Lieb--Seiringer de la conjecture de
Bessis--Moussa--Villani de 1975, réglée par Herbert Stahl en 2011 par un argument d'analyse complexe et
simplifiée par Eremenko en 2015 ; la tâche honnête consiste donc à reconstruire une démonstration ou à
trouver une voie plus courte relevant de la seule analyse matricielle.
\end{problem}

\noindent Bessis, Moussa et Villani demandèrent en 1975 si
$\operatorname{tr} e^{-(A + tB)}$ est la transformée de Laplace d'une mesure positive, ce qui aurait
justifié les approximants de Padé que les physiciens appliquaient déjà aux séries de perturbation en
mécanique statistique quantique. La reformulation de Lieb et Seiringer en 2004 a transformé cette
question analytique en l'énoncé fini et purement algébrique ci-dessus --- et il a encore tenu en échec
les analystes matriciels pendant sept ans, parce que développer $(A + tB)^m$ produit des mots en deux
lettres non commutantes dont les traces refusent d'être comparées terme à terme. La démonstration de
Stahl est venue d'une direction inattendue, la théorie des surfaces de Riemann et l'asymptotique des
polynômes orthogonaux, et en 2023 Heinävaara l'a déduite d'un théorème de structure pour les matrices
hermitiennes. C'est une bonne étude de cas sur la durée pendant laquelle une inégalité matricielle
\og{} manifestement vraie \fg{} peut résister.

\begin{refs}
\item Contexte : Théorème de Stahl --- Wikipédia (en anglais) (\url{https://en.wikipedia.org/wiki/Stahl%27s_theorem})
\item Article : H. R. Stahl, \og{} Proof of the BMV conjecture \fg{} (2013), Acta Mathematica
\item Lecture : R. Horn \& C. Johnson, \textit{Topics in Matrix Analysis}, ch. 6
\end{refs}

\bigskip


\section*{Pour aller plus loin}


\vfill
\noindent\emph{Hors de la Caverne} --- des probl\`emes, pas de r\'eponses.
\end{document}
